\documentclass[11pt, a4paper]{book}
\usepackage[utf8]{inputenc}
\usepackage[T1]{fontenc}
\usepackage[english]{babel}
\usepackage{amsmath}
\usepackage{amsfonts}
\usepackage{amssymb}
\usepackage{geometry}
\usepackage{titlesec}
\usepackage{enumitem}
\usepackage{xcolor}
\usepackage{hyperref}
\usepackage{fancyhdr}
\usepackage{tcolorbox}
\usepackage{tabularx}
\usepackage{booktabs}
\usepackage{parskip}
\usepackage{titletoc}
\usepackage{multicol}
\usepackage{blindtext}

\geometry{left=2.5cm, right=2.5cm, top=2.5cm, bottom=2.5cm}

\definecolor{primary}{RGB}{0, 51, 102}
\definecolor{secondary}{RGB}{0, 102, 153}
\definecolor{accent}{RGB}{200, 50, 50}
\definecolor{lightgray}{RGB}{245, 245, 245}
\definecolor{solbg}{RGB}{240, 255, 240}
\definecolor{a1color}{RGB}{34, 139, 34}
\definecolor{a2color}{RGB}{0, 102, 204}
\definecolor{b1color}{RGB}{204, 102, 0}
\definecolor{b2color}{RGB}{153, 0, 153}
\definecolor{c1color}{RGB}{200, 0, 0}
\definecolor{bluebox}{RGB}{230, 240, 255}
\definecolor{greenbox}{RGB}{230, 255, 240}
\definecolor{yellowbox}{RGB}{255, 250, 230}
\definecolor{redbox}{RGB}{255, 240, 240}

\hypersetup{
    colorlinks=true,
    linkcolor=primary,
    filecolor=primary,
    urlcolor=secondary,
    pdfpagelabels=true,
    pdfstartview=Fit,
}

\makeatletter
\newcommand\fmtcolorbox[1]{%
  \@ifnextchar[%]
    {\@fmtcolorbox{#1}}%
    {\@fmtcolorbox{#1}[o]}%
}
\def\@fmtcolorbox#1[#2]#3{%
  \newtcolorbox{mybox}[1]{colframe=#1,colback=#3,arc=0mm,outer arc=0mm,
    boxsep=0mm,left=3pt,right=3pt,top=3pt,bottom=3pt,before skip=0mm,after skip=0mm}%
  \begin{mybox}{}%
}
\makeatother

\newtcolorbox{infobox}[1]{colframe=primary,colback=bluebox,arc=3mm,rounded all,
    boxsep=5mm,left=5pt,right=5pt,top=5pt,bottom=5pt,
    fonttitle=\bfseries, title={#1}}

\newtcolorbox{warnbox}[1]{colframe=accent,colback=redbox,arc=3mm,rounded all,
    boxsep=5mm,left=5pt,right=5pt,top=5pt,bottom=5pt,
    fonttitle=\bfseries, title={#1}}

\newtcolorbox{tippbox}[1]{colframe=b1color,colback=yellowbox,arc=3mm,rounded all,
    boxsep=5mm,left=5pt,right=5pt,top=5pt,bottom=5pt,
    fonttitle=\bfseries, title={#1}}

\newtcolorbox{levelbox}[1]{colframe=a1color,colback=greenbox,arc=3mm,rounded all,
    boxsep=5mm,left=5pt,right=5pt,top=5pt,bottom=5pt,
    fonttitle=\bfseries, title={#1}}

\makeatletter
\newcommand{\chapterend}{\vfill\pagebreak}
\makeatother

\pagestyle{fancy}
\fancyhf{}
\rhead{\textit{Deutsche Grammatik A1-C}}
\lhead{George Mastro}
\rfoot{Seite \thepage}
\fancypagestyle{plain}{\rfoot{Seite \thepage}}

\title{\textbf{\Huge Deutsche Grammatik}\\ \Large Kompendium A1-C\\ \small Basierend auf Grammatik aktiv, Lehr- und Übungsbuch, B-Grammatik}
\author{George Mastro}
\date{April 2026}

\begin{document}

\pdfbookmark[1]{Titelseite}{title}
\maketitle
\thispagestyle{empty}
\vspace{2cm}

\begin{center}
\textbf{\LARGE Inhaltsverzeichnis mit Links}\\[1cm]
\tableofcontents
\end{center}

\newpage
\thispagestyle{empty}
\quad
\newpage

\frontmatter
\chapter{Einleitung}
Dieses Grammatikwerk bietet eine vollständige Referenz der deutschen Grammatik von A1 bis C1-Niveau. Die Inhalte wurden aus drei Hauptwerken zusammengestellt:

\section{Referenzwerke}
\begin{enumerate}
    \item \textbf{Grammatik aktiv B1+/B2/C1} - Übungsbuch mit theoretischen Erklärungen und praktischen Übungen
    \item \textbf{Lehr- und Übungsbuch der deutschen Grammatik} - Umfassendes Nachschlagewerk für Fortgeschrittene
    \item \textbf{B-Grammatik} - Didaktisch aufbereitete Grammatik für den Unterricht
\end{enumerate}

\section{Aufbau}
Das Werk ist nach grammatischen Kategorien organisiert:
\begin{itemize}
    \item \textbf{Teil I: Das Nomen} - Deklination, Artikel, N-Deklination
    \item \textbf{Teil II: Das Verb} - Zeitformen, Modalverben, Konjugation
    \item \textbf{Teil III: Syntax} - Satzbau, Nebensätze, Passiv
    \item \textbf{Teil IV: Präpositionen}
    \item \textbf{Teil V: Adjektive und Partikeln}
    \item \textbf{Teil VI: Ergänzungen}
    \item \textbf{Teil VII: Anhang}
\end{itemize}

\section{Legende}
\begin{itemize}
    \item \textcolor{a1color}{\textbf{A1}} - Grundstufe
    \item \textcolor{a2color}{\textbf{A2}} - Grundstufe II
    \item \textcolor{b1color}{\textbf{B1}} - Mittelstufe
    \item \textcolor{b2color}{\textbf{B2}} - Mittelstufe II
    \item \textcolor{c1color}{\textbf{C1}} - Oberstufe
\end{itemize}

\section{Navigationshilfen}
\begin{itemize}
    \item Klicken Sie auf Kapitelüberschriften im Inhaltsverzeichnis, um direkt zum Thema zu springen
    \item Im PDF: Verwenden Sie die Seitennavigation oder Lesezeichen
    \item Rückverweise: \hyperlink{cha:Deklination}{Zurück zum Anfang}
\end{itemize}

\mainmatter

\part{Das Nomen und seine Begleiter}
\pdfbookmark[0]{Teil I: Das Nomen}{part1}
\chapter{Deklination der Nomen}
\label{cha:Deklination}

Die Deklination im Deutschen markiert die grammatische Funktion eines Nomens durch Kasus (Nominativ, Akkusativ, Dativ, Genitiv). Im Gegensatz zu analytischen Sprachen wie dem Englischen behält das Deutsche eine starke flektierende Struktur.

\section{Grundlagen der Kasus}
\label{sec:Kasus}

Die vier Kasus im Deutschen:

\subsection{Nominativ (Wer? Was?)}
\begin{tcolorbox}[title=Funktion]
Der Nominativ bezeichnet das Subjekt des Satzes - wer oder was handelt oder existiert.
\end{tcolorbox}

\begin{itemize}
    \item \textbf{Der} Mann kommt.
    \item \textbf{Das} Kind schläft.
    \item \textbf{Die} Katze ist klein.
    \item \textbf{Ich} bin müde.
    \item \textbf{Sie} arbeitet.
    \item \textbf{Wir} sind fertig.
    \item \textbf{Die} Studenten lernen.
\end{itemize}

\subsection{Akkusativ (Wen? Was?)}
\begin{tcolorbox}[title=Funktion]
Der Akkusativ bezeichnet das direkte Objekt - wen oder was wird direkt betroffen.
\end{tcolorbox}

\begin{itemize}
    \item Ich sehe \textbf{den} Mann.
    \item Sie hat \textbf{das} Buch gelesen.
    \item Er kennt \textbf{die} Frau.
    \item Wir kaufen \textbf{die} Äpfel.
    \item Ich meine \textbf{dich}.
    \item Er liebt \textbf{seine} Frau.
    \item Wir brauchen \textbf{einen} Wagen.
\end{itemize}

\subsection{Dativ (Wem?)}
\begin{tcolorbox}[title=Funktion]
Der Dativ bezeichnet das indirekte Objekt - wem wird etwas gegeben oder zugeordnet.
\end{tcolorbox}

\begin{itemize}
    \item Ich helfe \textbf{dem} Mann.
    \item Sie gibt \textbf{der} Katze Futter.
    \item Er schreibt \textbf{dem} Freund.
    \item Wir danken \textbf{ihr}.
    \item Das gehört \textbf{mir}.
    \item Sie vertraut \textbf{ihm}.
    \item Ich antworte \textbf{dem} Lehrer.
\end{itemize}

\subsection{Genitiv (Wessen?)}
\begin{tcolorbox}[title=Funktion]
Der Genitiv bezeichnet Besitz oder Zugehörigkeit - wessen etwas ist.
\end{tcolorbox}

\begin{itemize}
    \item Das ist das Auto \textbf{des} Mannes.
    \item Die Farbe \textbf{der} Blume ist rot.
    \item Das Haus \textbf{des} Nachbarn ist groß.
    \item Die Mutter \textbf{des Kindes}.
    \item Das Buch \textbf{meines Bruders}.
    \item Die Tür \textbf{des Hauses}.
    \item Der Wagen \textbf{meines Vaters}.
\end{itemize}

\section{Deklinationstabellen}
\label{sec:Deklinationstabellen}

\subsection{Schwache Deklination (bestimmter Artikel)}
\begin{center}
\begin{tabular}{|l|c|c|c|c|}
\hline
\textbf{Kasus} & \textbf{Maskulin} & \textbf{Feminin} & \textbf{Neutral} & \textbf{Plural} \\ \hline
Nominativ & der & die & das & die \\ \hline
Akkusativ & den & die & das & die \\ \hline
Dativ & dem & der & dem & den (+n) \\ \hline
Genitiv & des & der & des & der \\ \hline
\end{tabular}
\end{center}

\subsection{Schwache Deklination (unbestimmter Artikel)}
\begin{center}
\begin{tabular}{|l|c|c|c|c|}
\hline
\textbf{Kasus} & \textbf{Maskulin} & \textbf{Feminin} & \textbf{Neutral} & \textbf{Plural} \\ \hline
Nominativ & ein & eine & ein & (keine) \\ \hline
Akkusativ & einen & eine & ein & (keine) \\ \hline
Dativ & einem & einer & einem & (keinen) \\ \hline
Genitiv & eines & einer & eines & (keiner) \\ \hline
\end{tabular}
\end{center}

\section{Übungen zur Deklination}
\label{sec:DeklinationUebungen}

\begin{tcolorbox}[title={Aufgabe}]
Bestimmen Sie den Kasus und bilden Sie den korrekten Artikel.
\end{tcolorbox}

\begin{enumerate}[label=\textbf{Aufgabe \arabic*}.]
    \item Ich warte auf \underline{\hspace{1.5cm}} Bus. \quad \textit{(Lösung: den -- Akkusativ)}
    \item Das Auto \underline{\hspace{1.5cm}} Nachbarn ist neu. \quad \textit{(Lösung: des -- Genitiv)}
    \item Sie hilft \underline{\hspace{1.5cm}} alten Mann. \quad \textit{(Lösung: dem -- Dativ)}
    \item \underline{\hspace{1.5cm}} Frau gab mir \underline{\hspace{1.5cm}} Buch. \quad \textit{(Lösung: Die, ein -- Nominativ, Akkusativ)}
    \item Wir denken an \underline{\hspace{1.5cm}} Zukunft. \quad \textit{(Lösung: die -- Akkusativ)}
    \item \underline{\hspace{1.5cm}} Kindern geht es gut. \quad \textit{(Lösung: Den -- Dativ)}
    \item Trotz \underline{\hspace{1.5cm}} Regens gingen wir los. \quad \textit{(Lösung: des -- Genitiv)}
    \item Ich habe \underline{\hspace{1.5cm}} Apfel gegessen. \quad \textit{(Lösung: einen -- Akkusativ)}
    \item Er spricht mit \underline{\hspace{1.5cm}} Lehrerin. \quad \textit{(Lösung: der -- Dativ)}
    \item Das ist das Haus \underline{\hspace{1.5cm}} Freundes. \quad \textit{(Lösung: des -- Genitiv)}
    \item Für \underline{\hspace{1.5cm}} gute Arbeit danke ich dir. \quad \textit{(Lösung: die -- Akkusativ)}
    \item \underline{\hspace{1.5cm}} Studenten lernen fleißig. \quad \textit{(Lösung: Die -- Nominativ)}
    \item Während \underline{\hspace{1.5cm}} Pause tranken wir Kaffee. \quad \textit{(Lösung: der -- Genitiv)}
    \item Ich gebe \underline{\hspace{1.5cm}} Mädchen ein Geschenk. \quad \textit{(Lösung: dem -- Dativ)}
    \item \underline{\hspace{1.5cm}} Berg ist hoch. \quad \textit{(Lösung: Der -- Nominativ)}
    \item Wir fahren nach \underline{\hspace{1.5cm}} Schweiz. \quad \textit{(Lösung: der -- Dativ/Akkusativ fix)}
    \item Ohne \underline{\hspace{1.5cm}} Geld kannst du nichts kaufen. \quad \textit{(Lösung: - -- Nullartikel)}
    \item \underline{\hspace{1.5cm}} Eltern sind stolz. \quad \textit{(Lösung: Die -- Nominativ)}
    \item Ich erinnere mich an \underline{\hspace{1.5cm}} Sommer. \quad \textit{(Lösung: den -- Akkusativ)}
    \item \underline{\hspace{1.5cm}} Probleme sind gelöst. \quad \textit{(Lösung: Die -- Nominativ)}
    \item Das ist der Wunsch \underline{\hspace{1.5cm}} Königs. \quad \textit{(Lösung: des -- Genitiv)}
    \item Der Rat \underline{\hspace{1.5cm}} Freundes war wichtig. \quad \textit{(Lösung: des -- Genitiv)}
    \item Ich gedenke \underline{\hspace{1.5cm}} Helden. \quad \textit{(Lösung: der -- Genitiv)}
    \item Er bedarf \underline{\hspace{1.5cm}} Hilfe. \quad \textit{(Lösung: der -- Genitiv)}
    \item Sie wurde \underline{\hspace{1.5cm}} Lage Herr. \quad \textit{(Lösung: der -- Genitiv)}
    \item Die Lösung \underline{\hspace{1.5cm}} Problems war einfach. \quad \textit{(Lösung: des -- Genitiv)}
    \item Der Erfolg \underline{\hspace{1.5cm}} Projekts freut uns. \quad \textit{(Lösung: des -- Genitiv)}
    \item Das liegt am Wetter \underline{\hspace{1.5cm}} Woche. \quad \textit{(Lösung: der -- Genitiv)}
    \item Der Chef \underline{\hspace{1.5cm}} Abteilung ist streng. \quad \textit{(Lösung: der -- Genitiv)}
    \item Die Höhe \underline{\hspace{1.5cm}} Berge ist impressive. \quad \textit{(Lösung: der -- Genitiv)}
\end{enumerate}

\section{Erweiterte Übungen: Kasusbestimmung}
\label{sec:ErweiterteUebungen}

Bestimmen Sie den Kasus der unterstrichenen Nomen:

\begin{enumerate}
    \item \underline{\hspace{1.5cm}} Junge hat Hunger. (Nominativ)
    \item Ich sehe \underline{\hspace{1.5cm}} junge Mann. (Akkusativ)
    \item Sie gibt \underline{\hspace{1.5cm}} Kind ein Buch. (Dativ)
    \item Das Auto \underline{\hspace{1.5cm}} Nachbarin ist rot. (Genitiv)
    \item Wir helfen \underline{\hspace{1.5cm}} alten Frau. (Dativ)
    \item Der Lehrer erklärt \underline{\hspace{1.5cm}} Schüler die Regel. (Dativ, Akkusativ)
    \item Ich erinnere mich an \underline{\hspace{1.5cm}} schöne Tag. (Akkusativ)
    \item Trotz \underline{\hspace{1.5cm}} Krankheit ging er zur Arbeit. (Genitiv)
    \item Der Preis \underline{\hspace{1.5cm}} Reise war hoch. (Genitiv)
    \item Sie ist stolz auf \underline{\hspace{1.5cm}} Erfolg. (Akkusativ)
    \item Das ist das Buch \underline{\hspace{1.5cm}} Autors. (Genitiv)
    \item Ich vertraue \underline{\hspace{1.5cm}} Freund. (Dativ)
    \item Der Brief \underline{\hspace{1.5cm}} Schwester liegt auf dem Tisch. (Genitiv)
    \item Sie freut sich über \underline{\hspace{1.5cm}} Geschenk. (Akkusativ)
    \item Wegen \underline{\hspace{1.5cm}} Staubs fiel der Unterricht aus. (Genitiv)
\end{enumerate}

\section{Regionale Besonderheiten}
\label{sec:RegionaleBesonderheiten}

\subsection{Genitiv-Ersatz durch von + Dativ}
In der gesprochenen Umgangssprache wird der Genitiv oft durch „von + Dativ" ersetzt:
\begin{itemize}
    \item \textit{Standard}: Das Auto des Mannes
    \item \textit{Umgangssprache}: Das Auto von dem Mann / vom Mann
    \item \textit{Standard}: Das Buch des Kindes
    \item \textit{Umgangssprache}: Das Buch von dem Kind
\end{itemize}

Dies gilt jedoch als grammatikalisch inkorrekt in der Standardsprache und im schriftlichen Gebrauch.

\section{Lösungen zu den Übungen}
\label{sec:Loesungen}

Die Lösungen zu allen Übungen finden Sie im Anhang.

\chapterend
\chapter{Artikel}
\label{cha:Artikel}

Artikel begleiten das Nomen und bestimmen seinen Kasus, Numerus und Genus.

\section{Bestimmter Artikel}
\label{sec:BestimmterArtikel}

\begin{center}
\begin{tabular}{|l|c|c|c|c|}
\hline
\textbf{Kasus} & \textbf{Maskulin} & \textbf{Feminin} & \textbf{Neutral} & \textbf{Plural} \\ \hline
Nominativ & der & die & das & die \\ \hline
Akkusativ & den & die & das & die \\ \hline
Dativ & dem & der & dem & den \\ \hline
Genitiv & des & der & des & der \\ \hline
\end{tabular}
\end{center}

\section{Unbestimmter Artikel}
\label{sec:UnbestimmterArtikel}

\begin{center}
\begin{tabular}{|l|c|c|c|c|}
\hline
\textbf{Kasus} & \textbf{Maskulin} & \textbf{Feminin} & \textbf{Neutral} & \textbf{Plural} \\ \hline
Nominativ & ein & eine & ein & keine \\ \hline
Akkusativ & einen & eine & ein & keine \\ \hline
Dativ & einem & einer & einem & keinen \\ \hline
Genitiv & eines & einer & eines & keiner \\ \hline
\end{tabular}
\end{center}

\section{Nullartikel}
\label{sec:Nullartikel}

Der Nullartikel wird verwendet bei:
\begin{itemize}
    \item \textbf{Abstrakta}: Glück, Freude, Liebe
    \item \textbf{Stoffnamen}: Wasser, Brot, Fleisch
    \item \textbf{Berufen/Nationalitäten}: Er ist Arzt. Sie ist Deutsche.
    \item \textbf{plurale Nomen ohne Artikel}: Kinder, Bücher
    \item \textbf{fixen Wendungen}: Hunger haben, Angst haben
\end{itemize}

\section{Gebrauch des Artikels}
\label{sec:ArtikelGebrauch}

\subsection{Bestimmt vs. Unbestimmt}
\begin{itemize}
    \item \textbf{Bestimmt}: Etwas Bekanntes, bereits Erwähntes, Einzigartiges
    \item \textbf{Unbestimmt}: Etwas Neues, Unbekanntes
\end{itemize}

\section{Übungen}
\label{sec:ArtikelUebungen}

\begin{enumerate}
    \item Ich brauche \underline{\hspace{1cm}} Buch. (ein)
    \item Das ist \underline{\hspace{1cm}} Haus. (das)
    \item \underline{\hspace{1cm}} Kind spielt im Garten. (ein)
    \item Ich kenne \underline{\hspace{1cm}} Mann nicht. (der)
    \item Sie hat \underline{\hspace{1cm}}Katze. (eine)
\end{enumerate}

\chapterend
\chapter{N-Deklination}
\label{cha:NDeklination}

Die N-Deklination betrifft eine spezifische Gruppe maskuliner Nomen, die in allen Fällen außer Nominativ Singular eine Endung erhalten.

\section{Nomen der N-Deklination}
\label{sec:NDeklinationNomen}

\subsection{Nomen auf -e}
\begin{itemize}
    \item der Junge $\rightarrow$ dem Jungen, den Jungen
    \item der Kunde $\rightarrow$ dem Kunden, den Kunden
    \item der Kollege $\rightarrow$ dem Kollegen, den Kollegen
    \item der Arbeiter $\rightarrow$ dem Arbeiter, die Arbeiter
\end{itemize}

\subsection{Nomen auf -ant, -ent, -ist}
\begin{itemize}
    \item der Praktikant $\rightarrow$ dem Praktikanten
    \item der Student $\rightarrow$ dem Studenten
    \item der Polizist $\rightarrow$ dem Polizisten
    \item der Journalist $\rightarrow$ dem Journalisten
\end{itemize}

\subsection{Ausnahmen (unregelmäßige N-Deklination)}
\begin{itemize}
    \item der Herr $\rightarrow$ dem Herrn, die Herren
    \item der Mensch $\rightarrow$ dem Menschen, die Menschen
    \item der Nachbar $\rightarrow$ dem Nachbarn, die Nachbarn
    \item der Bauer $\rightarrow$ dem Bauern, die Bauern
    \item der Held $\rightarrow$ dem Helden, die Helden
    \item der Prinz $\rightarrow$ dem Prinzen, die Prinzen
    \item der Kunde $\rightarrow$ dem Kunden, die Kunden
    \item der Experte $\rightarrow$ dem Experten, die Experten
    \item der Architekt $\rightarrow$ dem Architekten, die Architekten
\end{itemize}

\section{Deklinationstabelle}
\label{sec:NDeklinationTabelle}

\begin{center}
\begin{tabular}{|l|c|c|}
\hline
\textbf{Kasus} & \textbf{Singular} & \textbf{Plural} \\ \hline
Nominativ & der Architekt & die Architekten \\ \hline
Akkusativ & den Architekten & die Architekten \\ \hline
Dativ & dem Architekten & den Architekten \\ \hline
Genitiv & des Architekten & der Architekten \\ \hline
\end{tabular}
\end{center}

\section{Übungen zur N-Deklination}
\label{sec:NDeklinationUebungen}

\begin{tcolorbox}[title={Aufgabe}]
Deklinieren Sie die folgenden Nomen im Dativ Singular und Plural.
\end{tcolorbox}

\begin{enumerate}[label=\textbf{Aufgabe \arabic*}.]
    \item der Architekt (Dativ Sg.: \underline{\hspace{1cm}}, Pl.: \underline{\hspace{1cm}})
    \item der Kunde (Dativ Sg.: \underline{\hspace{1cm}}, Pl.: \underline{\hspace{1cm}})
    \item der Student (Dativ Sg.: \underline{\hspace{1cm}}, Pl.: \underline{\hspace{1cm}})
    \item der Herr (Dativ Sg.: \underline{\hspace{1cm}}, Pl.: \underline{\hspace{1cm}})
    \item der Mensch (Dativ Sg.: \underline{\hspace{1cm}}, Pl.: \underline{\hspace{1cm}})
\end{enumerate}

\chapterend
\chapter{Nomen-Verb-Verbindungen}
\label{cha:NomenVerb}

Viele Nomen kommen in festen Verbindungen mit Verben vor. Diese Verbindungen sind charakteristisch für den deutschen Sprachgebrauch.

\section{Verben mit Akkusativ}
\label{sec:VerbenMitAkk}

\begin{itemize}
    \item \textbf{abgeben}: Ich gebe den Job ab.
    \item \textbf{anfangen}: Wir fangen ein neues Projekt an.
    \item \textbf{anrufen}: Kannst du mich anrufen?
    \item \textbf{aufmachen}: Mach das Fenster auf!
    \item \textbf{ausfüllen}: Ich fülle das Formular aus.
    \item \textbf{besuchen}: Wir besuchen die Großeltern.
    \item \textbf{bezahlen}: Er bezahlt die Rechnung.
    \item \textbf{brauchen}: Ich brauche ein neues Auto.
    \item \textbf{durchsuchen}: Die Polizei durchsucht das Haus.
    \item \textbf{einladen}: Sie lädt ihre Freunde ein.
    \item \textbf{erzählen}: Er erzählt eine Geschichte.
    \item \textbf{finden}: Wir finden die Lösung.
    \item \textbf{fragen}: Ich frage den Lehrer.
    \item \textbf{haben}: Er hat ein Haus.
    \item \textbf{kennen}: Kennst du ihn?
    \item \textbf{kochen}: Sie kocht das Abendessen.
    \item \textbf{legen}: Leg das Buch auf den Tisch!
    \item \textbf{lesen}: Ich lese ein Buch.
    \item \textbf{lieben}: Er liebt seine Frau.
    \item \textbf{mitnehmen}: Nimm das Geschenk mit!
    \item \textbf{mögen}: Ich mag Kaffee.
    \item \textbf{nehmen}: Nimm den Bus!
    \item \textbf{rufen}: Rufen Sie einen Arzt!
    \item \textbf{schenken}: Sie schenkt mir Blumen.
    \item \textbf{schicken}: Ich schicke eine E-Mail.
    \item \textbf{schreiben}: Schreibe einen Brief!
    \item \textbf{sehen}: Ich sehe den Film.
    \item \textbf{setzen}: Setz dich hin!
    \item \textbf{stellen}: Stell die Flasche hin!
    \item \textbf{suchen}: Ich suche meine Brille.
    \item \textbf{tragen}: Er trägt eine Tasche.
    \item \textbf{treffen}: Wir treffen uns um 18 Uhr.
    \item \textbf{trinken}: Trinkst du Kaffee?
    \item \textbf{verstehen}: Verstehst du Deutsch?
    \item \textbf{waschen}: Ich wasche das Auto.
    \item \textbf{werfen}: Wirft du den Ball?
\end{itemize}

\section{Verben mit Dativ}
\label{sec:VerbenMitDat}

Verben, die ein Dativ-Objekt verlangen:

\subsection{Gefühlsverben}
\begin{itemize}
    \item \textbf{gefallen}: Das Buch gefällt mir.
    \item \textbf{gehören}: Das Auto gehört mir.
    \item \textbf{schmecken}: Der Kuchen schmeckt uns.
    \item \textbf{fehlen}: Mir fehlt nur noch ein Buch.
    \item \textbf{passen}: Diese Schuhe passen dir.
\end{itemize}

\subsection{Hilfsverben}
\begin{itemize}
    \item \textbf{helfen}: Ich helfe meiner Schwester.
    \item \textbf{danken}: Ich danke dir.
    \item \textbf{zustimmen}: Ich stimme dir zu.
    \item \textbf{glauben}: Ich glaube dir.
    \item \textbf{folgen}: Der Hund folgt seinem Herrn.
\end{itemize}

\subsection{Objektsverben}
\begin{itemize}
    \item \textbf{antworten}: Er antwortet mir.
    \item \textbf{begegnen}: Wir begegnen ihr.
    \item \textbf{gratulieren}: Ich gratuliere dir.
    \item \textbf{hören}: Ich höre dir zu.
    \item \textbf{zuhören}: Ich höre dir zu.
\end{itemize}

\subsection{Empfindungsverben}
\begin{itemize}
    \item \textbf{wehtun}: Der Arm tut mir weh.
    \item \textbf{leidtun}: Es tut mir leid.
\end{itemize}

\section{Verben mit Dativ und Akkusativ}
\label{sec:VerbenMitDatAkk}

Diese Verben verlangen zwei Objekte:

\begin{itemize}
    \item \textbf{bringen}: Ich bringe meiner Mutter Blumen.
    \item \textbf{empfehlen}: Ich empfehle dir diesen Film.
    \item \textbf{erklären}: Ich erkläre dem Studenten die Regel.
    \item \textbf{geben}: Ich gebe dir das Buch.
    \item \textbf{kaufen}: Ich kaufe meinem Vater eine Uhr.
    \item \textbf{kochen}: Ich koche meinem Mann Suppe.
    \item \textbf{schenken}: Ich schenke meiner Freundin einen Ring.
    \item \textbf{schicken}: Ich schicke dir eine Nachricht.
    \item \textbf{schreiben}: Ich schreibe meinem Chef einen Bericht.
    \item \textbf{stehlen}: Er stiehlt dem Museum ein Gemälde.
    \item \textbf{wünschen}: Ich wünsche dir alles Gute.
    \item \textbf{zeigen}: Ich zeige dem Touristen den Weg.
\end{itemize}

\section{Verben mit Genitiv}
\label{sec:VerbenMitGen}

Diese Verben sind im schriftlichen Deutsch, besonders akademisch, wichtig:

\begin{itemize}
    \item \textbf{bedürfen}: Der Patient bedarf medizinischer Hilfe.
    \item \textbf{gedenken}: Wir gedenken der Toten.
    \item \textbf{sich annehmen}: Sie nahm sich des Problems an.
    \item \textbf{sich bedienen}: Der Professor bedient sich klassischer Zitate.
    \item \textbf{sich erfreuen}: Das Projekt erfreut sich großer Beliebtheit.
    \item \textbf{Herr werden}: Er wurde der Lage Herr.
\end{itemize}

\section{Übungen}
\label{sec:NomenVerbUebungen}

\begin{enumerate}
    \item Kannst du \underline{\hspace{1cm}} helfen? (ich) \quad \textit{(Lösung: mir)}
    \item Das Kleid passt \underline{\hspace{1cm}} sehr gut. (du) \quad \textit{(Lösung: dir)}
    \item Wem glaubst du? (er) \quad \textit{(Lösung: ihm)}
    \item Ich gebe \underline{\hspace{1cm}} das Geschenk. (du) \quad \textit{(Lösung: dir)}
    \item Er zeigt \underline{\hspace{1cm}} das Foto. (wir) \quad \textit{(Lösung: uns)}
\end{enumerate}

\chapterend

\part{Das Verb}
\pdfbookmark[0]{Teil II: Das Verb}{part2}
\chapter{Zeitformen}
\label{cha:Zeitformen}

Die Zeitformen des deutschen Verbs ermöglichen die präzise zeitliche Einordnung von Handlungen.

\section{Präsens (Gegenwart)}
\label{sec:Praesens}

Das Präsens wird verwendet für:
\begin{itemize}
    \item \textbf{Gegenwart}: Ich arbeite heute.
    \item \textbf{Zukunft}: Morgen fliege ich nach Berlin.
    \item \textbf{historische Gegenwart}: 1789 bricht die Französische Revolution aus.
\end{itemize}

\section{Präteritum (Vergangenheit)}
\label{sec:Praeteritum}

Das Präteritum wird verwendet für:
\begin{itemize}
    \item \textbf{abgeschlossene Handlungen in der Vergangenheit}: Ich sah den Film gestern.
    \item \textbf{Erzählsprache}: Er ging nach Hause.
    \item \textbf{gewohnte Handlungen}: Als Kind ging er jeden Tag in die Schule.
\end{itemize}

\section{Perfekt (Vergangenheit)}
\label{sec:Perfekt}

Das Perfekt wird verwendet für:
\begin{itemize}
    \item \textbf{abgeschlossene Handlungen in der gesprochenen Sprache}: Ich habe den Film gesehen.
    \item \textbf{heute abgeschlossene Handlungen}: Heute habe ich viel gearbeitet.
\end{itemize}

\subsection{Perfekt mit Modalverben}
\label{sec:PerfektModal}

Bei Modalverben gibt es zwei Formen:

\begin{tcolorbox}[title=Standardform (Doppelter Infinitiv)]
\hfill\newline
haben/sein + Infinitiv (Modalverb) + Partizip II (Vollverb)
\end{tcolorbox}

\begin{itemize}
    \item Ich habe arbeiten müssen. (= Ich musste arbeiten.)
    \item Sie hat nicht kommen können. (= Sie konnte nicht kommen.)
    \item Er hat das Buch lesen wollen. (= Er wollte das Buch lesen.)
\end{itemize}

\begin{tcolorbox}[title=Ersatzform]
\hfill\newline
haben/sein + zu + Infinitiv + Partizip II (gehabt)
\end{tcolorbox}

\begin{itemize}
    \item Ich habe zu arbeiten gehabt. (= Ich musste arbeiten.)
    \item Wir haben früher zu gehen gehabt. (= Wir mussten früher gehen.)
\end{itemize}

\section{Plusquamperfekt (Vorvergangenheit)}
\label{sec:Plusquamperfekt}

Bildung: hatte/waren + Partizip II

\begin{itemize}
    \item Nachdem sie die Prüfung bestanden hatte, feierte sie.
    \item Als ich ankam, waren sie schon gegangen.
    \item Bevor er anrief, hatte ich ihn gesucht.
\end{itemize}

\section{Futur I (Zukunft)}
\label{sec:FuturI}

Bildung: werden + Infinitiv

\begin{itemize}
    \item \textbf{Zukunft}: Morgen werde ich arbeiten.
    \item \textbf{Vermutung}: Er wird etwa 40 Jahre alt sein.
    \item \textbf{Willen/Befehl}: Du wirst das tun!
\end{itemize}

\section{Futur II (Zukunft + Abgeschlossen)}
\label{sec:FuturII}

Bildung: werden + Partizip II + haben/sein

\begin{itemize}
    \item \textbf{abgeschlossene Zukunft}: Nächstes Jahr werde ich das Examen gemacht haben.
    \item \textbf{Vermutung über Vergangenheit}: Er wird schon gegangen sein.
\end{itemize}

\section{Übungen zu Zeitformen}
\label{sec:ZeitformenUebungen}

\subsection{Perfekt mit Modalverben}
\begin{enumerate}
    \item „Ich darf das nicht tun." $\rightarrow$ Perfekt \quad \textit{(Lösung: Ich habe das nicht tun dürfen.)}
    \item „Sie kann nicht kommen." $\rightarrow$ Perfekt \quad \textit{(Lösung: Sie hat nicht kommen können.)}
    \item „Er will das Buch lesen." $\rightarrow$ Perfekt \quad \textit{(Lösung: Er hat das Buch lesen wollen.)}
    \item „Wir müssen früher gehen." $\rightarrow$ Ersatzform \quad \textit{(Lösung: Wir haben früher zu gehen gehabt.)}
    \item „Ich soll die Hausaufgaben machen." $\rightarrow$ Perfekt \quad \textit{(Lösung: Ich habe die Hausaufgaben machen sollen.)}
    \item „Sie darf nicht rauchen." $\rightarrow$ Perfekt \quad \textit{(Lösung: Sie hat nicht rauchen dürfen.)}
    \item „Er kann gut schwimmen." $\rightarrow$ Perfekt \quad \textit{(Lösung: Er hat gut schwimmen können.)}
\end{enumerate}

\subsection{Plusquamperfekt}
\begin{enumerate}
    \item Nachdem sie die Prüfung \underline{\hspace{1cm}}, feierte sie. (bestehen) \quad \textit{(Lösung: hatte bestanden)}
    \item Als ich ankam, \underline{\hspace{1cm}} sie. (gehen) \quad \textit{(Lösung: war gegangen)}
    \item Bevor er anrief, \underline{\hspace{1cm}} ich ihn. (suchen) \quad \textit{(Lösung: hatte gesucht)}
    \item Sie war traurig, weil ihr Hund \underline{\hspace{1cm}}. (sterben) \quad \textit{(Lösung: war gestorben)}
    \item Ich \underline{\hspace{1cm}} den Zug, deshalb kam ich zu spät. (verpassen) \quad \textit{(Lösung: hatte verpasst)}
    \item Als du klingeltest, \underline{\hspace{1cm}} ich. (duschen) \quad \textit{(Lösung: hatte geduscht)}
    \item Wir \underline{\hspace{1cm}} noch nie in Berlin, bevor wir reisten. (sein) \quad \textit{(Lösung: waren gewesen)}
\end{enumerate}

\subsection{Futur}
\begin{enumerate}
    \item Morgen \underline{\hspace{1cm}} ich. (reisen) \quad \textit{(Lösung: werde reisen)}
    \item Das \underline{\hspace{1cm}} teuer. (sein) \quad \textit{(Lösung: wird teuer sein)}
    \item Nächstes Jahr \underline{\hspace{1cm}} ich das Examen. (machen) \quad \textit{(Lösung: werde das Examen gemacht haben)}
    \item Wenn du kommst, \underline{\hspace{1cm}} ich schon. (gehen) \quad \textit{(Lösung: werde ich schon gegangen sein)}
    \item \underline{\hspace{1cm}} du pünktlich? (sein) \quad \textit{(Lösung: Wirst du pünktlich sein?)}
    \item Um 20 Uhr \underline{\hspace{1cm}} wir essen. (fertig sein) \quad \textit{(Lösung: werden wir essen fertig sein)}
    \item \underline{\hspace{1cm}} sie das Buch schon? (lesen) \quad \textit{(Lösung: Wird sie das Buch schon gelesen haben?)}
\end{enumerate}

\chapterend
\chapter{Modalverben}
\label{cha:Modalverben}

Modalverben drücken Modalitäten wie Möglichkeit, Erlaubnis, Verpflichtung, Wunsch oder Vermutung aus.

\section{Übersicht der Modalverben}
\label{sec:ModalverbenUebersicht}

\begin{center}
\begin{tabular}{|l|c|c|c|c|}
\hline
\textbf{Verb} & \textbf{Bedeutung} & \textbf{Konjugation} & \textbf{Praeteritum} & \textbf{Perfekt} \\ \hline
dürfen & Erlaubnis, Möglichkeit & darf, darfst, darf, dürfen, dürft, dürfen & durfte & habe...dürfen \\ \hline
können & Fähigkeit, Möglichkeit & kann, kannst, kann, können, könnt, können & konnte & habe...können \\ \hline
müssen & Verpflichtung & muss, musst, muss, müssen, müsst, müssen & musste & habe...müssen \\ \hline
sollen & Empfehlung, Pflicht & soll, sollst, soll, sollen, sollt, sollten & sollte & habe...sollen \\ \hline
wollen & Wollen, Absicht & will, willst, will, wollen, wollt, wollen & wollte & habe...wollen \\ \hline
mögen & Sympathie, Wunsch & mag, magst, mag, mögen, mögt, mögen & mochte & habe...gemocht \\ \hline
\end{tabular}
\end{center}

\section{Grammatik aktiv: Modalverben im Satz}
\label{sec:ModalverbenSatz}

\begin{enumerate}
    \item Ich \underline{\hspace{1cm}} Deutsch lernen. (wollen) \quad \textit{(Lösung: will)}
    \item Du \underline{\hspace{1cm}} das nicht tun. (dürfen) \quad \textit{(Lösung: darfst)}
    \item Er \underline{\hspace{1cm}} gut schwimmen. (können) \quad \textit{(Lösung: kann)}
    \item Wir \underline{\hspace{1cm}} heute arbeiten. (müssen) \quad \textit{(Lösung: müssen)}
    \item Ihr \underline{\hspace{1cm}} pünktlich sein. (sollen) \quad \textit{(Lösung: sollt)}
    \item Sie \underline{\hspace{1cm}} Kaffee. (mögen) \quad \textit{(Lösung: mögen)}
\end{enumerate}

\section{Subjektive Bedeutung der Modalverben}
\label{sec:SubjektiveModalverben}

Modalverben können nicht nur objektive Modalitäten ausdrücken, sondern auch subjektive Einschätzungen:

\subsection{Weitergabe von Informationen (sollen/wollen)}
\begin{itemize}
    \item \textbf{sollen}: Man gibt Informationen wieder, die man gehört/gelesen hat. (Der Wahrheitsgehalt ist unsicher.)
        \begin{itemize}
            \item Er soll der beste Spieler sein. (= Man sagt, er sei der beste...)
            \item Sie soll in Paris leben. (= Ich habe gehört, dass sie in Paris lebt.)
        \end{itemize}
    \item \textbf{wollen}: Jemand behauptet etwas über sich selbst. (Distanzierung des Sprechers.)
        \begin{itemize}
            \item Er will der Beste sein. (= Er behauptet von sich, der Beste zu sein.)
            \item Sie will alles gewusst haben. (= Sie behauptet, alles gewusst zu haben.)
        \end{itemize}
\end{itemize}

\subsection{Ausdruck einer Vermutung}
\begin{itemize}
    \item \textbf{müssen/müsste}: Hohe Wahrscheinlichkeit (90-99\%)
        \begin{itemize}
            \item Er muss noch im Büro sein. (= Ich bin sicher, dass er noch im Büro ist.)
            \item Das muss stimmen. (= Es ist logisch, dass es stimmt.)
        \end{itemize}
    \item \textbf{dürfte}: Wahrscheinlichkeit (75\%)
        \begin{itemize}
            \item Er dürfte noch im Büro sein. (= Wahrscheinlich ist er noch im Büro.)
            \item Das dürfte falsch sein. (= Es ist wahrscheinlich falsch.)
        \end{itemize}
    \item \textbf{kann/könnte}: Möglichkeit (50\%)
        \begin{itemize}
            \item Das kann sein. (= Es ist möglich.)
            \item Das könnte stimmen. (= Es könnte vielleicht stimmen.)
        \end{itemize}
\end{itemize}

\section{Übungen: Subjektive Modalverben}
\label{sec:SubjektiveModalUebungen}

\begin{enumerate}
    \item Er \underline{\hspace{1cm}} der beste Spieler sein. (sollen/wollen - Gerücht) \quad \textit{(Lösung: soll)}
    \item Er \underline{\hspace{1cm}} der Beste sein. (wollen - Eigenbehauptung) \quad \textit{(Lösung: will)}
    \item Das \underline{\hspace{1cm}} stimmen. (müssen - logischer Schluss) \quad \textit{(Lösung: muss)}
    \item Er \underline{\hspace{1cm}} noch im Büro sein. (dürfte - wahrscheinlich) \quad \textit{(Lösung: dürfte)}
    \item Das \underline{\hspace{1cm}} falsch sein. (könnte - möglich) \quad \textit{(Lösung: könnte)}
    \item Sie \underline{\hspace{1cm}} krank sein. (sollen - Information) \quad \textit{(Lösung: soll)}
    \item Er \underline{\hspace{1cm}} das nicht gewusst haben. (wollen - Distanzierung) \quad \textit{(Lösung: will)}
\end{enumerate}

\section{Konjunktiv II Formen der Modalverben}
\label{sec:ModalKonjunktivII}

\begin{center}
\begin{tabular}{|l|c|c|}
\hline
\textbf{Infinitiv} & \textbf{Konjunktiv II} & \textbf{Bedeutung} \\ \hline
dürfen & dürfte & dürfte (höflich) \\ \hline
könnte & könnte & könnte (höflich) \\ \hline
müsste & müsste & müsste (höflich) \\ \hline
sollte & sollte & sollte (höflich) \\ \hline
wollte & wollte & wollte (unhöflich, nur „möchten" verwenden) \\ \hline
möchte & möchte & möchte (höflich) \\ \hline
\end{tabular}
\end{center}

\chapterend
\chapter{Reflexive Verben}
\label{cha:Reflexiv}

Reflexive Verben werden mit einem Reflexivpronomen gebraucht. Das Reflexivpronomen zeigt an, dass sich die Handlung auf das Subjekt des Satzes bezieht.

\section{Reflexivpronomen}
\label{sec:Reflexivpronomen}

\begin{center}
\begin{tabular}{|l|c|c|}
\hline
\textbf{Person} & \textbf{Akkusativ} & \textbf{Dativ} \\ \hline
ich & mich & mir \\ \hline
du & dich & dir \\ \hline
er/sie/es & sich & sich \\ \hline
wir & uns & uns \\ \hline
ihr & euch & euch \\ \hline
sie/Sie & sich & sich \\ \hline
\end{tabular}
\end{center}

\section{Arten von Reflexivverben}
\label{sec:ReflexivArten}

\subsection{Echte Reflexivverben}
Das Reflexivpronomen ist obligatorisch:
\begin{itemize}
    \item Ich wasche \textbf{mich}.
    \item Er setzt \textbf{sich}.
    \item Wir freuen \textbf{uns}.
    \item Sie entscheidet \textbf{sich}.
\end{itemize}

\subsection{Reflexiv mit Dativ}
Das reflexive Objekt steht im Dativ:
\begin{itemize}
    \item Ich wasche \textbf{mir} die Hände.
    \item Er kocht \textbf{sich} ein Ei.
    \item Sie nimmt \textbf{sich} Zeit.
\end{itemize}

\subsection{Reziproke Verben (wechselseitig)}
\begin{itemize}
    \item Wir sehen \textbf{einander}.
    \item Sie helfen \textbf{einander}.
    \item Sie verstehen \textbf{einander}.
\end{itemize}

\section{Häufige reflexive Verben}
\label{sec:ReflexivVerben}

\begin{itemize}
    \item \textbf{sich anmelden}: Ich melde mich an.
    \item \textbf{sich anziehen}: Er zieht sich an.
    \item \textbf{sich aufregen}: Sie regt sich auf.
    \item \textbf{sich beeilen}: Beeil dich!
    \item \textbf{sich beschweren}: Ich beschwere mich.
    \item \textbf{sich entscheiden}: Ich entscheide mich.
    \item \textbf{sich entschließen}: Er hat sich entschlossen.
    \item \textbf{sich erholen}: Wir erholen uns.
    \item \textbf{sich erkälten}: Ich habe mich erkältet.
    \item \textbf{sich erinnern}: Ich erinnere mich.
    \item \textbf{sich freuen}: Ich freue mich.
    \item \textbf{sich fühlen}: Ich fühle mich wohl.
    \item \textbf{sich interessieren}: Ich interessiere mich für Kunst.
    \item \textbf{sich konzentrieren}: Ich konzentriere mich.
    \item \textbf{sich langweilen}: Er langweilt sich.
    \item \textbf{sich leisten}: Das kann ich mir nicht leisten.
    \item \textbf{sich patschen}: Das Kind patscht sich.
    \item \textbf{sich streiten}: Sie streiten sich.
    \item \textbf{sich unterhalten}: Wir unterhalten uns.
    \item \textbf{sich verlieben}: Er hat sich verliebt.
    \item \textbf{sich verlaufen}: Ich habe mich verlaufen.
    \item \textbf{sich vorstellen}: Ich stelle mir etwas vor.
\end{itemize}

\section{Übungen zu reflexiven Verben}
\label{sec:ReflexivUebungen}

\begin{enumerate}
    \item Ich beschwere \underline{\hspace{1cm}} über die neue Waschmaschine. (mich)
    \item Die Gäste begrüßen \underline{\hspace{1cm}}. (sich)
    \item Ich sehe \underline{\hspace{1cm}} den Film nicht noch einmal an. (mir)
    \item Er interessiert \underline{\hspace{1cm}} für Kunst. (sich)
    \item Wir freuen \underline{\hspace{1cm}} auf den Urlaub. (uns)
    \item Du musst \underline{\hspace{1cm}} beeilen. (dich)
    \item Sie erinnert \underline{\hspace{1cm}} an den Termin. (sich)
    \item Ich kann \underline{\hspace{1cm}} das nicht leisten. (mir)
    \item Er bildet \underline{\hspace{1cm}} viel ein. (sich)
    \item Wir unterhalten \underline{\hspace{1cm}} gut. (uns)
\end{enumerate}

\section{Grammatik aktiv: Besonderheiten}
\label{sec:ReflexivBesonderheiten}

\begin{tcolorbox}[title=Wichtig]
Bei manchen Verben ändert sich die Bedeutung mit/reflexiv:
\end{tcolorbox}

\begin{itemize}
    \item \textbf{erinnern}: Ich erinnere ihn an den Termin. (= Ich sage ihm, dass...)
    \item \textbf{sich erinnern}: Ich erinnere mich an den Termin. (= Ich habe das Gedächtnis daran.)
    \item \textbf{treffen}: Ich treffe ihn. (= Wir begegnen uns.)
    \item \textbf{sich treffen}: Wir treffen uns. (= Wir haben eine Verabredung.)
\end{itemize}

\chapterend
\chapter{Konjunktiv}
\label{cha:Konjunktiv}

Der Konjunktiv ist ein Modus, der Irrealität, Wünsche, Höflichkeit oder indirekte Rede ausdrückt.

\section{Konjunktiv I – Indirekte Rede}
\label{sec:KonjunktivI}

Der Konjunktiv I wird verwendet, um fremde Aussagen neutral wiederzugeben.

\subsection{Bildung}
\begin{center}
\begin{tabular}{|l|c|c|c|}
\hline
\textbf{Person} & \textbf{sein} & \textbf{haben} & \textbf{werden} \\ \hline
ich & sei & habe & werde \\ \hline
du & seiest & habest & werdest \\ \hline
er/sie/es & sei & habe & werde \\ \hline
wir & seien & haben & werden \\ \hline
ihr & seiet & habet & werdet \\ \hline
sie/Sie & seien & haben & werden \\ \hline
\end{tabular}
\end{center}

\subsection{Verwendung}
\begin{itemize}
    \item \textbf{Indirekte Rede}: Er sagt, er \textbf{sei} krank. (Original: „Ich bin krank.")
    \item \textbf{Zeitungsberichte}: Der Minister \textbf{teilte} mit, die Verhandlungen \textbf{seien} erfolgreich.
    \item \textbf{Wissenschaftliche Berichte}: Die Studie \textbf{zeige}, dass...
\end{itemize}

\section{Konjunktiv II – Irreale Wünsche und Höflichkeit}
\label{sec:KonjunktivII}

Der Konjunktiv II drückt Irrealität, Wünsche oder Höflichkeit aus.

\subsection{Bildung}
\begin{itemize}
    \item \textbf{regelmäßige Verben}: Präteritum + -e
        \begin{itemize}
            \item kommen $\rightarrow$ käme
            \item gehen $\rightarrow$ ginge
            \item sehen $\rightarrow$ sähe
        \end{itemize}
    \item \textbf{unregelmäßige Verben}: Umlaut im Präteritum
        \begin{itemize}
            \item haben $\rightarrow$ hätte
            \item sein $\rightarrow$ wäre
            \item werden $\rightarrow$ würde
        \end{itemize}
    \item \textbf{Modalverben}: Siehe Tabelle
        \begin{itemize}
            \item können $\rightarrow$ könnte
            \item müssen $\rightarrow$ müsste
            \item dürfen $\rightarrow$ dürfte
        \end{itemize}
\end{itemize}

\subsection{Höflichkeitsform}
\begin{itemize}
    \item \textbf{Könnten Sie mir helfen?} (= Bitte helfen Sie mir.)
    \item \textbf{Würden Sie das bitte wiederholen?} (= Bitte wiederholen Sie das.)
    \item \textbf{Ich würde gern wissen, ob...} (= Ich möchte gern wissen, ob...)
\end{itemize}

\subsection{Ireale Wünsche}
\begin{itemize}
    \item \textbf{Wäre ich nur reich!} (= Ich bin nicht reich, aber ich wünschte, ich wäre es.)
    \item \textbf{Wenn ich nur Zeit hätte!} (= Ich habe keine Zeit.)
    \item \textbf{Es wäre schön, wenn...} (= Es ist nicht so.)
\end{itemize}

\subsection{Ireale Bedingungssätze}
\begin{itemize}
    \item \textbf{Wenn ich Zeit hätte, würde ich kommen.}
    \item \textbf{Wäre ich früher gekommen, hätte ich ihn getroffen.}
\end{itemize}

\section{Übungen zum Konjunktiv}
\label{sec:KonjunktivUebungen}

\subsection{Konjunktiv I}
\begin{enumerate}
    \item „Ich komme morgen." $\rightarrow$ indirekt \quad \textit{(Lösung: Er sagt, er komme morgen.)}
    \item „Wir haben Hunger." $\rightarrow$ indirekt \quad \textit{(Lösung: Sie sagen, sie hätten Hunger.)}
    \item „Ich kann schwimmen." $\rightarrow$ indirekt \quad \textit{(Lösung: Er sagt, er könne schwimmen.)}
    \item „Es ist kalt." $\rightarrow$ indirekt \quad \textit{(Lösung: Er sagt, es sei kalt.)}
    \item „Ich bin fertig." $\rightarrow$ indirekt \quad \textit{(Lösung: Sie sagt, sie sei fertig.)}
    \item „Wir werden warten." $\rightarrow$ indirekt \quad \textit{(Lösung: Sie sagen, sie würden warten.)}
    \item „Ich lese gern." $\rightarrow$ indirekt \quad \textit{(Lösung: Er sagt, er lese gern.)}
    \item „Die Aufgabe war leicht." $\rightarrow$ indirekt \quad \textit{(Lösung: Sie sagen, die Aufgabe sei leicht gewesen.)}
    \item „Ich habe gelernt." $\rightarrow$ indirekt \quad \textit{(Lösung: Er sagt, er habe gelernt.)}
    \item „Wir müssen arbeiten." $\rightarrow$ indirekt \quad \textit{(Lösung: Sie sagen, sie müssten arbeiten.)}
\end{enumerate}

\subsection{Konjunktiv II}
\begin{enumerate}
    \item Ich möchte wissen, ob du Zeit hast. $\rightarrow$ höflich \quad \textit{(Lösung: Ich würde gern wissen, ob du Zeit hättest.)}
    \item Ich bin nicht reich. $\rightarrow$ Wunsch \quad \textit{(Lösung: Ich wäre gern reich.)}
    \item Kannst du mir helfen? $\rightarrow$ höflicher \quad \textit{(Lösung: Könntest du mir helfen?)}
    \item Ich habe keine Zeit. $\rightarrow$ Wunsch \quad \textit{(Lösung: Ich hätte gern Zeit.)}
    \item Bin ich du? $\rightarrow$ Irreal \quad \textit{(Lösung: Wäre ich du,...)}
    \item Ich möchte Kaffee. $\rightarrow$ höflich \quad \textit{(Lösung: Hätten Sie gern Kaffee?)}
    \item Wenn ich Zeit \underline{\hspace{1cm}}, \underline{\hspace{1cm}} ich. (haben/reisen) \quad \textit{(Lösung: hätte - würde reisen)}
    \item Ich \underline{\hspace{1cm}} gern in Berlin. (sein) \quad \textit{(Lösung: wäre)}
    \item \underline{\hspace{1cm}} du mir sagen, wie spät es ist? (können) \quad \textit{(Lösung: Könntest)}
    \item Ich \underline{\hspace{1cm}} gern ein Auto. (haben) \quad \textit{(Lösung: hätte)}
\end{enumerate}

\chapterend
\chapter{Funktionsverbgefüge (FVG)}
\label{cha:FVG}

Funktionsverbgefüge sind Verbindungen aus einem Nomen und einem „leeren" Verb (haben, sein, bringen, stellen), die ein einfaches Verb ersetzen. Sie sind charakteristisch für den formellen, akademischen Stil.

\section{Häufige FVG}
\label{sec:FVGHaeufig}

\begin{center}
\begin{tabular}{|l|l|}
\hline
\textbf{einfaches Verb} & \textbf{FVG} \\ \hline
beeinflussen & Einfluss nehmen auf \\ \hline
kritisieren & Kritik üben an \\ \hline
sich entscheiden & eine Entscheidung treffen \\ \hline
helfen & Hilfe leisten \\ \hline
fragen & eine Frage stellen \\ \hline
beginnen & den Anfang machen mit \\ \hline
enden & ein Ende setzen \\ \hline
diskutieren & zur Diskussion stehen \\ \hline
erlauben & die Erlaubnis geben \\ \hline
beabsichtigen & die Absicht haben \\ \hline
sich bewerben & sich bewerben um \\ \hline
versuchen & den Versuch machen \\ \hline
beweisen & den Beweis erbringen \\ \hline
erklären & eine Erklärung abgeben \\ \hline
verstehen & Verständnis haben für \\ \hline
\end{tabular}
\end{center}

\section{Übungen zu FVG}
\label{sec:FVGUebungen}

\begin{enumerate}
    \item Wir diskutieren die Pläne. $\rightarrow$ \textit{(Lösung: Die Pläne stehen zur Diskussion.)}
    \item Der Arzt kritisiert das Verhalten. $\rightarrow$ \textit{(Lösung: Der Arzt übt Kritik am Verhalten.)}
    \item Sie beginnen mit der Präsentation. $\rightarrow$ \textit{(Lösung: Sie machen den Anfang mit der Präsentation.)}
    \item Er beweist seine Theorie. $\rightarrow$ \textit{(Lösung: Er erbringt den Beweis für seine Theorie.)}
    \item Wir treffen eine Entscheidung. $\rightarrow$ \textit{(Lösung: Wir treffen die Entscheidung.)}
\end{enumerate}

\section{Stilistische Hinweise}
\label{sec:FVGStil}

\begin{tcolorbox}[title=Stil]
\begin{itemize}
    \item \textbf{FVG}: formell, akademisch, schriftlich
    \item \textbf{einfaches Verb}: informell, mündlich
\end{itemize}
\end{tcolorbox}

FVG sollten im akademischen Schreiben verwendet werden, in der gesprochenen Sprache klingen sie jedoch oft steif.

\chapterend

\part{Satzbau und Syntax}
\pdfbookmark[0]{Teil III: Syntax}{part3}
\chapter{Satzbau}
\label{cha:Satzbau}

Der Satzbau im Deutschen folgt klaren Regeln, die die Wortstellung im Hauptsatz und Nebensatz bestimmen.

\section{Die Satzklammer}
\label{sec:Satzklammer}

Das deutsche Verb bildet eine „Satzklammer":

\begin{tcolorbox}[title=Satzklammer]
\begin{center}
Position 2: finites Verb \\
Ende: infinite Verbformen (Partizip, Infinitiv)
\end{center}
\end{tcolorbox}

\begin{itemize}
    \item Ich \textbf{sehe} den Film heute Abend \textbf{anzusehen}. (trennbares Verb)
    \item Er \textbf{hat} yesterday \textbf{arbeiten müssen}. (Modalverb)
    \item Wir \textbf{wollen} morgen ins Kino \textbf{gehen}.
\end{itemize}

\section{Hauptsatz-Wortstellung}
\label{sec:Hauptsatz}

\begin{center}
\begin{tabular}{|c|c|c|c|c|c|}
\hline
\textbf{Position 1} & \textbf{Position 2} & \textbf{Position 3} & \textbf{Position 4} & \textbf{Position 5} & \textbf{Position 6+} \\ \hline
Subjekt & finites Verb & Objekt 1 & Adverb & Objekt 2 & Verb 2... \\ \hline
\end{tabular}
\end{center}

Beispiele:
\begin{itemize}
    \item Ich \textbf{sehe} \underline{morgen} \underline{die} \underline{Freunde}.
    \item Heute \textbf{geht} \underline{er} \underline{nach} \underline{Hause}.
    \item \underline{Ich} \textbf{bin} \underline{nach} \underline{Frankfurt} \underline{gefahren}.
\end{itemize}

\section{Nebensatz-Wortstellung}
\label{sec:Nebensatz}

Im Nebensatz steht das Verb am Ende:

\begin{tcolorbox}[title=Nebensatz]
\begin{center}
Subjekt - Objekt - Adverb - \textbf{Verb am Ende}
\end{center}
\end{tcolorbox}

\begin{itemize}
    \item ...weil er \textbf{morgen} \underline{nach} \underline{Hause} \textbf{geht}.
    \item ...dass sie \underline{viel} \underline{Geld} \textbf{verdient}.
    \item ...obwohl es \textbf{kalt} \textbf{war}.
\end{itemize}

\section{Trennbare Verben}
\label{sec:Trennbar}

Im Hauptsatz: Finite Verbform Position 2, Partikel am Ende.
Im Nebensatz: Zusammen am Ende.

\begin{enumerate}
    \item Wann \underline{\hspace{1cm}} ihr in Berlin \underline{\hspace{1cm}}? (ankommen) \quad \textit{(Lösung: kommt ... an)}
    \item Um 8 Uhr \underline{\hspace{1cm}} der Unterricht \underline{\hspace{1cm}}. (anfangen) \quad \textit{(Lösung: fängt ... an)}
    \item Er \underline{\hspace{1cm}} mit dem Rauchen \underline{\hspace{1cm}}. (aufhören) \quad \textit{(Lösung: hört ... auf)}
    \item Wann \underline{\hspace{1cm}} du \underline{\hspace{1cm}}? (zurückkommen) \quad \textit{(Lösung: kommst ... zurück)}
    \item Ich \underline{\hspace{1cm}} ein Geschenk \underline{\hspace{1cm}}. (mitbringen) \quad \textit{(Lösung: bringe ... mit)}
    \item Das Kind \underline{\hspace{1cm}} um 20 Uhr \underline{\hspace{1cm}}. (einschlafen) \quad \textit{(Lösung: schläft ... ein)}
    \item Wann \underline{\hspace{1cm}} ihr \underline{\hspace{1cm}}? (vorbeikommen) \quad \textit{(Lösung: kommt ... vorbei)}
    \item Sie \underline{\hspace{1cm}} nach der Arbeit \underline{\hspace{1cm}}. (weggehen) \quad \textit{(Lösung: geht ... weg)}
    \item Ich \underline{\hspace{1cm}} dich später \underline{\hspace{1cm}}. (anrufen) \quad \textit{(Lösung: rufe ... an)}
    \item Kannst du das Fenster \underline{\hspace{1cm}}? (aufmachen) \quad \textit{(Lösung: aufmachen)}
\end{enumerate}

\chapterend
\chapter{Nebensätze}
\label{cha:Nebensaetze}

Nebensätze ergänzen oder erweitern den Hauptsatz. Sie werden durch Konjunktionen eingeleitet.

\section{Konjunktionen mit Verb am Ende}
\label{sec:Konjunktionen}

\subsection{Kausal (Ursache)}
\begin{itemize}
    \item \textbf{weil}: ..., weil sie krank war. (Nebensatz)
    \item \textbf{da}: Da ich müde war, ging ich ins Bett.
    \item \textbf{denn}: ..., denn ich war müde. (Hauptsatz!)
\end{itemize}

\subsection{Konzessiv (Gegensatz)}
\begin{itemize}
    \item \textbf{obwohl}: ..., obwohl es kalt war.
    \item \textbf{trotzdem}: ..., sie arbeitete trotzdem weiter. (Adverb, Hauptsatz!)
\end{itemize}

\subsection{Temporal (Zeit)}
\begin{itemize}
    \item \textbf{bevor}: ..., bevor ich ging.
    \item \textbf{nachdem}: ..., nachdem ich gegessen hatte.
    \item \textbf{als}: ..., als ich Kind war.
    \item \textbf{während}: ..., während ich las.
    \item \textbf{sobald}: ..., sobald ich kann.
    \item \textbf{bis}: ..., bis ich fertig bin.
\end{itemize}

\subsection{Konditional (Bedingung)}
\begin{itemize}
    \item \textbf{wenn/falls}: ..., wenn ich Zeit habe.
    \item \textbf{ohne dass}: ..., ohne dass ich wusste.
    \item \textbf{um zu}: ..., um zu lernen.
\end{itemize}

\subsection{Final (Zweck)}
\begin{itemize}
    \item \textbf{damit}: ..., damit du verstehst.
    \item \textbf{um zu}: ..., um zu lernen.
\end{itemize}

\section{Übungen zu Konjunktionen}
\label{sec:KonjunktionenUebungen}

\subsection{Kausal}
\begin{enumerate}
    \item Sie hat den Bus verpasst. Sie kam zu spät. $\rightarrow$ weil \quad \textit{(Lösung: Sie kam zu spät, weil sie den Bus verpasst hatte.)}
    \item Ich bleibe zu Hause. Ich bin krank. $\rightarrow$ denn \quad \textit{(Lösung: Ich bin krank, denn ich bleibe zu Hause.)}
    \item Es ist kalt. Wir bleiben drinnen. $\rightarrow$ da \quad \textit{(Lösung: Da es kalt ist, bleiben wir drinnen.)}
    \item Er hat nicht gelernt. Er hat die Prüfung nicht bestanden. $\rightarrow$ weil \quad \textit{(Lösung: Er hat die Prüfung nicht bestanden, weil er nicht gelernt hatte.)}
\end{enumerate}

\subsection{Konzessiv}
\begin{enumerate}
    \item Es war kalt. Wir sind schwimmen gegangen. $\rightarrow$ obwohl \quad \textit{(Lösung: Obwohl es kalt war, sind wir schwimmen gegangen.)}
    \item Er ist arm, aber er ist glücklich. $\rightarrow$ obwohl \quad \textit{(Lösung: Obwohl er arm ist, ist er glücklich.)}
    \item Sie hatte Kopfschmerzen. Sie arbeitete weiter. $\rightarrow$ trotzdem \quad \textit{(Lösung: Sie hatte Kopfschmerzen; sie arbeitete trotzdem weiter.)}
    \item Der Film war langweilig. Wir blieben bis zum Ende. $\rightarrow$ trotzdem \quad \textit{(Lösung: Der Film war langweilig; wir blieben trotzdem bis zum Ende.)}
\end{enumerate}

\chapterend
\chapter{Passiv}
\label{cha:Passiv}

Das Passiv drückt aus, dass die Handlung im Vordergrund steht und der Handelnde unbekannt oder unwichtig ist.

\section{Vorgangspassiv}
\label{sec:Vorgangspassiv}

Bildung: werden + Partizip II

\begin{itemize}
    \item Das Haus wird gebaut. (gerade)
    \item Die Rechnung wurde bezahlt. (vergangen)
    \item Das Buch wird gelesen werden. (Zukunft)
\end{itemize}

\section{Zustandspassiv}
\label{sec:Zustandspassiv}

Bildung: sein + Partizip II

\begin{itemize}
    \item Das Fenster ist geöffnet. (= Zustand)
    \item Die Tür war geschlossen.
    \item Die Arbeit ist gemacht.
\end{itemize}

\section{Passiversatzformen}
\label{sec:Passiversatz}

\subsection{man-Konstruktion}
\begin{itemize}
    \item Das Fahrrad wird repariert. $\rightarrow$ Man repariert das Fahrrad.
    \item Hier wird Deutsch gesprochen. $\rightarrow$ Man spricht hier Deutsch.
\end{itemize}

\subsection{sein + zu + Infinitiv}
\begin{itemize}
    \item Das Haus wird gebaut. $\rightarrow$ Das Haus ist zu bauen. (= man muss)
    \item Die Aufgabe ist zu lösen. (= man muss)
\end{itemize}

\subsection{lassen + sich}
\begin{itemize}
    \item Das Problem wird gelöst. $\rightarrow$ Das Problem lässt sich lösen.
    \item Der Text ist zu verstehen. $\rightarrow$ Der Text lässt sich verstehen.
\end{itemize}

\section{Übungen zum Passiv}
\label{sec:PassivUebungen}

\begin{enumerate}
    \item Jemand repariert das Fahrrad. $\rightarrow$ Passiv \quad \textit{(Lösung: Das Fahrrad wird repariert.)}
    \item Sie baut ein Haus. $\rightarrow$ Passiv \quad \textit{(Lösung: Ein Haus wird gebaut.)}
    \item Man hat das Fenster geöffnet. $\rightarrow$ Zustandspassiv \quad \textit{(Lösung: Das Fenster ist geöffnet.)}
    \item Jemand hat die Tür geschlossen. $\rightarrow$ Zustandspassiv \quad \textit{(Lösung: Die Tür ist geschlossen.)}
    \item Das Fahrrad wird repariert. $\rightarrow$ mit „man" \quad \textit{(Lösung: Man repariert das Fahrrad.)}
    \item Das Haus wird gebaut. $\rightarrow$ mit „sein+zu" \quad \textit{(Lösung: Das Haus ist zu bauen.)}
    \item Ich bringe das Auto zur Werkstatt. $\rightarrow$ mit „lassen" \quad \textit{(Lösung: Ich lasse das Auto reparieren.)}
\end{enumerate}

\chapterend
\chapter{Relativsätze}
\label{cha:Relativ}

Relativsätze werden durch Relativpronomen eingeleitet und beschreiben ein Nomen näher.

\section{Relativpronomen}
\label{sec:Relativpronomen}

Die Relativpronomen richten sich nach Genus, Numerus und Kasus des Bezugsworts:

\begin{center}
\begin{tabular}{|l|c|c|c|c|}
\hline
\textbf{Kasus} & \textbf{Maskulin} & \textbf{Feminin} & \textbf{Neutral} & \textbf{Plural} \\ \hline
Nominativ & der & die & das & die \\ \hline
Akkusativ & den & die & das & die \\ \hline
Dativ & dem & der & dem & denen \\ \hline
Genitiv & dessen & deren & dessen & deren \\ \hline
\end{tabular}
\end{center}

\section{Relativpronomen mit Präpositionen}
\label{sec:RelativPraep}

\begin{center}
\begin{tabular}{|l|l|}
\hline
\textbf{Präposition} & \textbf{Beispiel} \\ \hline
an + dem & an dem = an dem \\ \hline
in + dem & in dem = in dem \\ \hline
auf + das & auf das = auf das \\ \hline
mit + dem & mit dem = mit dem \\ \hline
über + das & über das = über das \\ \hline
\end{tabular}
\end{center}

\section{Übungen zu Relativsätzen}
\label{sec:RelativUebungen}

\subsection{Relativpronomen (Nom/Akk/Dat)}
\begin{enumerate}
    \item Das ist das Buch, \underline{\hspace{1cm}} mich interessiert. \quad \textit{(Lösung: das)}
    \item Der Mann, \underline{\hspace{1cm}} dort steht, ist mein Vater. \quad \textit{(Lösung: der)}
    \item Die Frau, \underline{\hspace{1cm}} singt, ist berühmt. \quad \textit{(Lösung: die)}
    \item Das ist der Mann, \underline{\hspace{1cm}} ich kenne. \quad \textit{(Lösung: den)}
    \item Die Frau, \underline{\hspace{1cm}} ich sah, war berühmt. \quad \textit{(Lösung: die)}
    \item Der Lehrer, \underline{\hspace{1cm}} ich Fragen stellte, antwortete. \quad \textit{(Lösung: dem)}
    \item Die Nachbarn, \underline{\hspace{1cm}} ich grüße, sind freundlich. \quad \textit{(Lösung: denen)}
    \item Das Land, \underline{\hspace{1cm}} ich entstamme, ist Griechenland. \quad \textit{(Lösung: dem)}
\end{enumerate}

\subsection{Relativsätze mit Präpositionen}
\begin{enumerate}
    \item Ich habe ein Problem. Ich spreche mit dem Lehrer. $\rightarrow$ Relativsatz \quad \textit{(Lösung: Ich habe ein Problem, mit dem ich mit dem Lehrer spreche.)}
    \item Sie arbeitet in einer Firma. Sie ist stolz darauf. $\rightarrow$ Relativsatz \quad \textit{(Lösung: Sie arbeitet in einer Firma, auf die sie stolz ist.)}
    \item Er hat einen Freund. Er vertraut ihm. $\rightarrow$ Relativsatz \quad \textit{(Lösung: Er hat einen Freund, dem er vertraut.)}
\end{enumerate}

\chapterend
\chapter{Nominalstil}
\label{cha:Nominalstil}

Der Nominalstil (auch „Nominalisierung") ersetzt Verben durch Substantive. Er ist charakteristisch für den formellen, akademischen Stil.

\section{Verben $\rightarrow$ Nomen}
\label{sec:VerbZuNomen}

\begin{center}
\begin{tabular}{|l|l|}
\hline
\textbf{Verb} & \textbf{Nomen} \\ \hline
entscheiden & die Entscheidung \\ \hline
planen & die Planung \\ \hline
analysieren & die Analyse \\ \hline
verbessern & die Verbesserung \\ \hline
entwickeln & die Entwicklung \\ \hline
erforschen & die Erforschung \\ \hline
bearbeiten & die Bearbeitung \\ \hline
beantworten & die Beantwortung \\ \hline
\end{tabular}
\end{center}

\section{Übungen zum Nominalstil}
\label{sec:NominalstilUebungen}

\begin{enumerate}
    \item Wir planen die Reise. $\rightarrow$ Nominalstil \quad \textit{(Lösung: Die Planung der Reise)}
    \item Sie analysieren die Daten. $\rightarrow$ Nominalstil \quad \textit{(Lösung: Die Analyse der Daten)}
    \item Er verbessert die Software. $\rightarrow$ Nominalstil \quad \textit{(Lösung: Die Verbesserung der Software)}
    \item Wir entwickeln neue Technologien. $\rightarrow$ Nominalstil \quad \textit{(Lösung: Die Entwicklung neuer Technologien)}
    \item Er forscht auf diesem Gebiet. $\rightarrow$ Nominalstil \quad \textit{(Lösung: Die Erforschung auf diesem Gebiet)}
\end{enumerate}

\section{Stilistische Hinweise}
\label{sec:NominalstilHinweise}

\begin{tcolorbox}[title=Stil]
\begin{itemize}
    \item \textbf{Nominalstil}: formell, akademisch, schriftlich
    \item \textbf{Verbalstil}: informell, mündlich
\end{itemize}
\end{tcolorbox}

Der Nominalstil sollte im akademischen Schreiben verwendet werden, in der gesprochenen Sprache klingt er jedoch oft steif.

\chapterend

\part{Präpositionen}
\pdfbookmark[0]{Teil IV: Präpositionen}{part4}
\chapter{Präpositionen}
\label{cha:Praepositionen}

Präpositionen verbinden Wörter und zeigen Beziehungen zwischen ihnen.

\section{Wechselpräpositionen (Dual Präpositionen)}
\label{sec:Wechselpraep}

Die Wechselpräpositionen können sowohl mit Dativ (Ort: Wo?) als auch mit Akkusativ (Richtung: Wohin?) verwendet werden:

\begin{center}
\begin{tabular}{|l|c|c|}
\hline
\textbf{Präposition} & \textbf{Wo? (Dativ)} & \textbf{Wohin? (Akkusativ)} \\ \hline
in & in der Schule & in die Schule \\ \hline
an & an der Wand & an die Wand \\ \hline
auf & auf dem Tisch & auf den Tisch \\ \hline
unter & unter dem Tisch & unter den Tisch \\ \hline
vor & vor dem Haus & vor das Haus \\ \hline
hinter & hinter dem Haus & hinter das Haus \\ \hline
neben & neben dem Haus & neben das Haus \\ \hline
zwischen & zwischen den Bäumen & zwischen die Bäume \\ \hline
über & über dem Tisch & über den Tisch \\ \hline
\end{tabular}
\end{center}

\begin{tcolorbox}[title=Wichtig]
\begin{itemize}
    \item \textbf{Ort (Stationär)}: Dativ
    \item \textbf{Richtung (Bewegung)}: Akkusativ
\end{itemize}
\end{tcolorbox}

\section{Übungen zu Wechselpräpositionen}
\label{sec:WechselpraepUebungen}

\begin{enumerate}
    \item Ich stelle die Blumen auf \underline{\hspace{1cm}} Tisch. (Richtung) \quad \textit{(Lösung: den)}
    \item Die Blumen stehen auf \underline{\hspace{1cm}} Tisch. (Ort) \quad \textit{(Lösung: dem)}
    \item Geh in \underline{\hspace{1cm}} Garten! (Richtung) \quad \textit{(Lösung: den)}
    \item Er ist in \underline{\hspace{1cm}} Garten. (Ort) \quad \textit{(Lösung: dem)}
    \item Häng das Bild an \underline{\hspace{1cm}} Wand. (Richtung) \quad \textit{(Lösung: die)}
    \item Das Bild hängt an \underline{\hspace{1cm}} Wand. (Ort) \quad \textit{(Lösung: der)}
    \item Setz dich vor \underline{\hspace{1cm}} Fernseher. (Richtung) \quad \textit{(Lösung: den)}
    \item Ich sitze vor \underline{\hspace{1cm}} Fernseher. (Ort) \quad \textit{(Lösung: dem)}
    \item Leg das Handy unter \underline{\hspace{1cm}} Kissen. (Richtung) \quad \textit{(Lösung: das)}
    \item Das Handy liegt unter \underline{\hspace{1cm}} Kissen. (Ort) \quad \textit{(Lösung: dem)}
    \item Wir laufen zwischen \underline{\hspace{1cm}} Bäume. (Richtung) \quad \textit{(Lösung: die)}
    \item Die Straße verläuft zwischen \underline{\hspace{1cm}} Bäumen. (Ort) \quad \textit{(Lösung: den)}
    \item Steck den Brief in \underline{\hspace{1cm}} Umschlag. (Richtung) \quad \textit{(Lösung: den)}
    \item Der Brief ist in \underline{\hspace{1cm}} Umschlag. (Ort) \quad \textit{(Lösung: dem)}
    \item Stell den Stuhl neben \underline{\hspace{1cm}} Sofa. (Richtung) \quad \textit{(Lösung: das)}
    \item Der Stuhl steht neben \underline{\hspace{1cm}} Sofa. (Ort) \quad \textit{(Lösung: dem)}
    \item Wir gehen über \underline{\hspace{1cm}} Brücke. (Richtung) \quad \textit{(Lösung: die)}
    \item Die Enten schwimmen unter \underline{\hspace{1cm}} Brücke. (Ort) \quad \textit{(Lösung: der)}
    \item Leg das Kleid über \underline{\hspace{1cm}} Stuhl. (Richtung) \quad \textit{(Lösung: den)}
    \item Das Kleid hängt über \underline{\hspace{1cm}} Stuhl. (Ort) \quad \textit{(Lösung: dem)}
\end{enumerate}

\section{Reine Dativ-Präpositionen}
\label{sec:DativPraep}

\begin{center}
\begin{tabular}{|l|l|}
\hline
\textbf{Präposition} & \textbf{Beispiel} \\ \hline
aus & Er kommt aus Berlin. \\ \hline
bei & Ich bin bei Maria. \\ \hline
mit & Ich fahre mit dem Bus. \\ \hline
nach & Nach dem Essen... \\ \hline
seit & Seit zwei Jahren... \\ \hline
von & Das Buch ist von ihm. \\ \hline
zu & Ich gehe zur Schule. \\ \hline
\end{tabular}
\end{center}

\section{Reine Akkusativ-Präpositionen}
\label{sec:AkkusativPraep}

\begin{center}
\begin{tabular}{|l|l|}
\hline
\textbf{Präposition} & \textbf{Beispiel} \\ \hline
bis & Bis morgen! \\ \hline
durch & Durch den Wald \\ \hline
gegen & Gegen die Wand \\ \hline
ohne & Ohne dich \\ \hline
um & Um die Ecke \\ \hline
\end{tabular}
\end{center}

\section{Genitiv-Präpositionen}
\label{sec:GenitivPraep}

\begin{center}
\begin{tabular}{|l|l|}
\hline
\textbf{Präposition} & \textbf{Beispiel} \\ \hline
trotz & Trotz des Regens... \\ \hline
wegen & Wegen des Verkehrs... \\ \hline
während & Während des Films... \\ \hline
statt & Statt des Bruders... \\ \hline
außerhalb & Außerhalb der Stadt \\ \hline
innerhalb & Innerhalb der Woche \\ \hline
aufgrund & Aufgrund der Situation... \\ \hline
oberhalb & Oberhalb des Berges \\ \hline
unterhalb & Unterhalb des Tales \\ \hline
diesseits & Diesseits des Flusses \\ \hline
jenseits & Jenseits des Meeres \\ \hline
\end{tabular}
\end{center}

\section{Übungen zu Genitiv-Präpositionen}
\label{sec:GenitivPraepUebungen}

\begin{enumerate}
    \item \underline{\hspace{1cm}} schlechten Wetters blieben wir zu Hause. (trotz) \quad \textit{(Lösung: Trotz des)}
    \item \underline{\hspace{1cm}} seiner Krankheit fehlte er. (wegen) \quad \textit{(Lösung: Wegen)}
    \item \underline{\hspace{1cm}} des Films schlief ich ein. (während) \quad \textit{(Lösung: Während)}
    \item \underline{\hspace{1cm}} des Verkehrs kam er zu spät. (wegen) \quad \textit{(Lösung: Wegen)}
    \item \underline{\hspace{1cm}} des Kaffees trinke ich Tee. (anstatt) \quad \textit{(Lösung: Anstatt)}
    \item \underline{\hspace{1cm}} seiner Krankheit blieb er zu Hause. (aufgrund) \quad \textit{(Lösung: Aufgrund)}
    \item \underline{\hspace{1cm}} des Parks ist es verboten. (innerhalb) \quad \textit{(Lösung: Innerhalb)}
    \item \underline{\hspace{1cm}} der Öffnungszeiten ist geschlossen. (außerhalb) \quad \textit{(Lösung: Außerhalb)}
    \item \underline{\hspace{1cm}} der Baumgrenze wächst nichts. (oberhalb) \quad \textit{(Lösung: Oberhalb)}
    \item \underline{\hspace{1cm}} des Meeres ist Europa. (diesseits) \quad \textit{(Lösung: Diesseits)}
    \item \underline{\hspace{1cm}} der Berge liegt ein Tal. (jenseits) \quad \textit{(Lösung: Jenseits)}
\end{enumerate}

\section{Feste Präpositionen mit festem Kasus}
\label{sec:FestePraep}

\begin{tcolorbox}[title={Präposition: Dativ}]
aus, bei, mit, nach, seit, von, zu
\end{tcolorbox}

\begin{tcolorbox}[title={Präposition: Akkusativ}]
bis, durch, gegen, ohne, um
\end{tcolorbox}

\begin{tcolorbox}[title={Präposition: Genitiv}]
trotz, wegen, während, statt
\end{tcolorbox}

\chapterend

\part{Adjektive und Partikeln}
\pdfbookmark[0]{Teil V: Adjektive}{part5}
\chapter{Adjektive}
\label{cha:Adjektive}

Adjektive beschreiben oder bestimmen Nomen näher.

\section{Deklination der Adjektive}
\label{sec:AdjDeklination}

Adjektive werden dekliniert und richten sich nach Genus, Numerus und Kasus des bezogenen Nomens.

\subsection{Starke Deklination (ohne Artikel)}
\begin{center}
\begin{tabular}{|l|c|c|c|c|}
\hline
\textbf{Kasus} & \textbf{Maskulin} & \textbf{Feminin} & \textbf{Neutral} & \textbf{Plural} \\ \hline
Nominativ & gut-er & gut-e & gut-es & gut-e \\ \hline
Akkusativ & gut-en & gut-e & gut-es & gut-e \\ \hline
Dativ & gut-em & gut-er & gut-em & gut-en \\ \hline
Genitiv & gut-en & gut-er & gut-en & gut-er \\ \hline
\end{tabular}
\end{center}

\subsection{Gemischte Deklination (unbestimmter Artikel)}
\begin{center}
\begin{tabular}{|l|c|c|c|c|}
\hline
\textbf{Kasus} & \textbf{Maskulin} & \textbf{Feminin} & \textbf{Neutral} & \textbf{Plural} \\ \hline
Nominativ & gut-er & gut-e & gut-es & gut-en \\ \hline
Akkusativ & gut-en & gut-e & gut-es & gut-en \\ \hline
Dativ & gut-en & gut-er & gut-en & gut-en \\ \hline
Genitiv & gut-en & gut-er & gut-en & gut-en \\ \hline
\end{tabular}
\end{center}

\subsection{Schwache Deklination (bestimmter Artikel)}
\begin{center}
\begin{tabular}{|l|c|c|c|c|}
\hline
\textbf{Kasus} & \textbf{Maskulin} & \textbf{Feminin} & \textbf{Neutral} & \textbf{Plural} \\ \hline
Nominativ & gut-e & gut-e & gut-e & gut-en \\ \hline
Akkusativ & gut-en & gut-e & gut-e & gut-en \\ \hline
Dativ & gut-en & gut-en & gut-en & gut-en \\ \hline
Genitiv & gut-en & gut-en & gut-en & gut-en \\ \hline
\end{tabular}
\end{center}

\section{Steigerung (Komparativ, Superlativ)}
\label{sec:AdjSteigerung}

\subsection{Regelmäßige Steigerung}
\begin{itemize}
    \item \textbf{Positiv}: klein
    \item \textbf{Komparativ}: kleiner
    \item \textbf{Superlativ}: am kleinsten
\end{itemize}

\subsection{Unregelmäßige Steigerung}
\begin{center}
\begin{tabular}{|l|c|c|}
\hline
\textbf{Positiv} & \textbf{Komparativ} & \textbf{Superlativ} \\ \hline
gut & besser & am besten \\ \hline
viel & mehr & am meisten \\ \hline
gern & lieber & am liebsten \\ \hline
groß & größer & am größten \\ \hline
hoch & höher & am höchsten \\ \hline
nah & näher & am nächsten \\ \hline
\end{tabular}
\end{center}

\section{Adjektiv vs. Adverb}
\label{sec:AdjvsAdv}

\begin{tcolorbox}[title=Unterscheidung]
\begin{itemize}
    \item \textbf{Adjektiv}: beschreibt das Nomen
    \item \textbf{Adverb}: beschreibt das Verb
\end{itemize}
\end{tcolorbox}

\begin{itemize}
    \item Er ist ein \textbf{guter} Schüler. (Adjektiv)
    \item Er lernt \textbf{gut}. (Adverb)
\end{itemize}

\section{Übungen zu Adjektiven}
\label{sec:AdjUebungen}

\begin{enumerate}
    \item Das ist ein \underline{\hspace{1cm}} Buch. (interessant)
    \item Er hat ein \underline{\hspace{1cm}} Auto. (neu)
    \item Die \underline{\hspace{1cm}} Frau ist meine Mutter. (alt)
    \item Ich habe \underline{\hspace{1cm}} Probleme. (viel)
    \item Das ist der \underline{\hspace{1cm}} Weg. (kurz)
\end{enumerate}

\chapterend
\chapter{Partikeln}
\label{cha:Partikeln}

Partikeln sind kleine Wörter, die die Stimmung, Einstellung oder Nuancen des Sprechers ausdrücken. Sie sind besonders wichtig für die mündliche Kommunikation.

\section{Abtönungspartikeln (Modalpartikeln)}
\label{sec:Modalpartikeln}

Abtönungspartikeln verändern die Bedeutung eines Satzes subtil:

\begin{center}
\begin{tabular}{|l|c|l|}
\hline
\textbf{Partikel} & \textbf{Funktion} & \textbf{Beispiel} \\ \hline
aber & Erstaunen, Bewertung & Das ist aber schön! \\ \hline
also & Schlussfolgerung & Also, das war's. \\ \hline
auch & Einbeziehung & Das kann ich auch. \\ \hline
bloß & Einschränkung & Beeil dich bloß! \\ \hline
denn & Frage & Wie geht's denn? \\ \hline
doch & Widerspruch/Erstaunen & Das ist doch klar! \\ \hline
eben & Resignation & Das ist eben so. \\ \hline
eigentlich & Einschränkung & Was meinst du eigentlich? \\ \hline
etwa & Frage & Ist das etwa wahr? \\ \hline
ja & Zustimmung & Das ist ja klar. \\ \hline
mal & Abschwächung & Komm mal vorbei! \\ \hline
nur & Einschränkung & Nur zu! \\ \hline
schon & Beruhigung & Das ist schon okay. \\ \hline
toll & Begeisterung & Das ist ja toll! \\ \hline
vielleicht & Erstaunen & Das ist vielleicht schwer! \\ \hline
wohl & Vermutung & Das wird wohl stimmen. \\ \hline
\end{tabular}
\end{center}

\section{Gradpartikeln}
\label{sec:Gradpartikeln}

Gradpartikeln verstärken oder relativieren die Bedeutung:

\begin{center}
\begin{tabular}{|l|c|l|}
\hline
\textbf{Partikel} & \textbf{Funktion} & \textbf{Beispiel} \\ \hline
sehr & Verstärkung & Das ist sehr gut. \\ \hline
ziemlich & Annäherung & Das ist ziemlich gut. \\ \hline
recht & Annäherung & Das ist recht gut. \\ \hline
ganz & Annäherung & Das ist ganz gut. \\ \hline
ein bisschen & Abschwächung & Ein bisschen besser. \\ \hline
ein wenig & Abschwächung & Ein wenig besser. \\ \hline
bisschen & Abschwächung & Das ist bisschen schwer. \\ \hline
\end{tabular}
\end{center}

\section{Steigerungspartikeln}
\label{sec:Steigerungspartikeln}

\begin{center}
\begin{tabular}{|l|c|l|}
\hline
\textbf{Partikel} & \textbf{Funktion} & \textbf{Beispiel} \\ \hline
noch & Steigerung & Das ist noch besser. \\ \hline
schon & Vorzeitigkeit & Das war schon gut. \\ \hline
immer & Wiederholung & Immer besser. \\ \hline
\end{tabular}
\end{center}

\section{negationspartikeln}
\label{sec:Negationspartikeln}

\begin{center}
\begin{tabular}{|l|c|l|}
\hline
\textbf{Partikel} & \textbf{Funktion} & \textbf{Beispiel} \\ \hline
nicht & Negation & Das ist nicht gut. \\ \hline
kein & Negation & Das ist kein Problem. \\ \hline
nie & Negation & Das war nie einfach. \\ \hline
niemals & starke Negation & Das ist niemals einfach. \\ \hline
keineswegs & starke Negation & Das ist keineswegs einfach. \\ \hline
\end{tabular}
\end{center}

\section{Antwortpartikeln}
\label{sec:Antwortpartikeln}

\begin{center}
\begin{tabular}{|l|c|l|}
\hline
\textbf{Partikel} & \textbf{Funktion} & \textbf{Beispiel} \\ \hline
ja & Zustimmung & Ja, das stimmt. \\ \hline
nein & Ablehnung & Nein, das stimmt nicht. \\ \hline
doch & Widerspruch & Nein, das ist doch falsch! \\ \hline
doch & positive Antwort & Hast du Zeit? - Doch. \\ \hline
vielleicht & Unsicherheit & Vielleicht. \\ \hline
na ja & skeptisch & Na ja, maybe. \\ \hline
\end{tabular}
\end{center}

\section{Übungen zu Partikeln}
\label{sec:PartikelnUebungen}

\begin{enumerate}
    \item Komm \underline{\hspace{1cm}} vorbei! (mal)
    \item Das ist \underline{\hspace{1cm}} schön! (aber)
    \item Wie geht's \underline{\hspace{1cm}}? (denn)
    \item Das ist \underline{\hspace{1cm}} klar! (doch)
    \item Beeil dich \underline{\hspace{1cm}}! (bloß)
\end{enumerate}

\chapterend

\part{Ergänzungen}
\pdfbookmark[0]{Teil VI: Ergänzungen}{part6}
\chapter{Wortstellung}
\label{cha:Wortstellung}

Die Wortstellung im Deutschen folgt bestimmten Regeln, die über die Grundpositionen hinausgehen.

\section{Stellung des Subjekts}
\label{sec:SubjektStellung}

\begin{itemize}
    \item \textbf{Position 1}: Normal: Subjekt
    \item \textbf{Position 1}: Topikalisierung: Andere Elemente
        \begin{itemize}
            \item Heute \textbf{geht} er einkaufen.
            \item Im Winter \textbf{schneit} es oft.
            \item Mir \textbf{geht} es gut.
        \end{itemize}
\end{itemize}

\section{Stellung der Negation}
\label{sec:NegationStellung}

\begin{itemize}
    \item \textbf{nicht}: Vor dem Verb oder am Ende
        \begin{itemize}
            \item Er \textbf{kommt nicht}.
            \item Er \textbf{hat} nicht \textbf{kommen} können.
        \end{itemize}
    \item \textbf{kein}: Vor dem Nomen
        \begin{itemize}
            \item Ich habe \textbf{kein} Geld.
        \end{itemize}
\end{itemize}

\section{Stellung von Adverbien}
\label{sec:AdverbStellung}

\begin{enumerate}
    \item Temporal (Wann?): oft am Anfang oder vor dem Verb
        \begin{itemize}
            \item \textbf{Heute} komme ich. / Ich komme \textbf{heute}.
        \end{itemize}
    \item Lokal (Wo?): nach dem Verb
        \begin{itemize}
            \item Ich bin \textbf{in Berlin}.
        \end{itemize}
    \item Modal (Wie?): vor dem Verb
        \begin{itemize}
            \item Er spricht \textbf{schnell}.
        \end{itemize}
\end{enumerate}

\section{Inversion (Umstellung)}
\label{sec:Inversion}

\begin{itemize}
    \item \textbf{Nach „und", „oder", „aber"}: Subjekt folgt
        \begin{itemize}
            \item Und \textbf{er} kam nach Hause.
            \item Aber \textbf{ich} blieb.
        \end{itemize}
    \item \textbf{Nach „da" (weil)}: Subjekt folgt
        \begin{itemize}
            \item Da \textbf{es} regnete, blieben wir.
        \end{itemize}
    \item \textbf{Nach „nicht nur...sondern auch"}: Subjekt folgt
        \begin{itemize}
            \item Nicht nur \textbf{ich}, sondern auch \textbf{er} kam.
        \end{itemize}
\end{itemize}

\chapterend
\chapter{Häufige Fehler und Stolpersteine}
\label{cha:Stolpersteine}

Dieses Kapitel behandelt die häufigsten Fehler und Stolpersteine der deutschen Grammatik.

\section{Häufige Kasus-Fehler}
\label{sec:KasusFehler}

\begin{enumerate}
    \item \textbf{Seit + Dativ}: Ich lerne seit zwei Jahren. (nicht: Jahr)
    \item \textbf{Anrufen + Akk}: Kannst du mich anrufen? (nicht: mir)
    \item \textbf{glauben + Dativ}: Ich glaube dir. (nicht: dir glauben)
    \item \textbf{helfen + Dativ}: Ich helfe dir. (nicht: dir helfen)
    \item \textbf{antworten + Dativ}: Ich antworte ihm. (nicht: auf ihn antworten - archaish)
\end{enumerate}

\section{Häufige Präpositionen-Fehler}
\label{sec:PraepFehler}

\begin{enumerate}
    \item \textbf{denken an + Akk}: Ich denke an dich. (nicht: dir)
    \item \textbf{warten auf + Akk}: Ich warte auf dich. (nicht: dir)
    \item \textbf{ab + Dativ}: Ich komme ab morgen. (nicht: von morgen)
    \item \textbf{außer + Dativ}: Außer mir war niemand da. (nicht: mich)
\end{enumerate}

\section{Häufige Verb-Fehler}
\label{sec:VerbFehler}

\begin{enumerate}
    \item \textbf{legen/stellen/setzen}:
        \begin{itemize}
            \item Ich lege das Buch auf den Tisch. (horizontal)
            \item Ich stelle das Buch auf den Tisch. (vertikal)
            \item Ich setze mich auf den Stuhl. (mit Bewegung)
        \end{itemize}
    \item \textbf{sein/werden}:
        \begin{itemize}
            \item Er ist müde. (Zustand)
            \item Er wird müde. (Veränderung)
        \end{itemize}
    \item \textbf{haben/sein (Hilfsverben)}:
        \begin{itemize}
            \item Ich habe gegessen. (transitiv)
            \item Ich bin gegangen. (intransitiv/bewegung)
        \end{itemize}
\end{enumerate}

\section{Häufige Artikel-Fehler}
\label{sec:ArtikelFehler}

\begin{enumerate}
    \item \textbf{Nullartikel}: Ich habe Hunger. (nicht: einen Hunger)
    \item \textbf{Nationalitäten}: Er ist Deutscher. (nicht: ein Deutscher - als Substantiv)
    \item \textbf{Berufe}: Sie ist Ärztin. (nicht: eine Ärztin - offiziell)
\end{enumerate}

\section{Häufige Wortstellung-Fehler}
\label{sec:WortstellungFehler}

\begin{enumerate}
    \item \textbf{Deshalb + Inversion}: ..., \textbf{deshalb bin ich} nicht gekommen. (nicht: deshalb ich bin nicht gekommen)
    \item \textbf{Wegen + Inversion}: Wegen des Wetters \textbf{bin ich} geblieben.
    \item \textbf{Relativpronomen}: Das ist der Mann, den ich gesehen habe. (nicht: welcher)
\end{enumerate}

\section{Häufige Zeiten-Fehler}
\label{sec:ZeitenFehler}

\begin{enumerate}
    \item \textbf{Perfekt vs. Präteritum}: In der gesprochenen Sprache wird Perfekt bevorzugt.
    \item \textbf{Futur I vs. Präsens}: Für Zukunft wird oft Präsens verwendet: Morgen fliege ich.
    \item \textbf{Plusquamperfekt}: Nur in Verbindung mit Perfekt/Präteritum: Nachdem ich gegessen hatte, ging ich.
\end{enumerate}

\section{Übersicht: Wichtige Unterscheidungen}
\label{sec:Unterscheidungen}

\begin{center}
\begin{tabular}{|l|l|l|}
\hline
\textbf{A} & \textbf{B} & \textbf{Beispiel} \\ \hline
kennen & wissen & Ich kenne ihn. / Ich weiß, dass... \\ \hline
legen & liegen & Ich lege das Buch. / Das Buch liegt. \\ \hline
stellen & stehen & Ich stelle das Buch. / Das Buch steht. \\ \hline
setzen & sitzen & Ich setze mich. / Ich sitze. \\ \hline
hängen & hängen & Ich hänge das Bild. / Das Bild hängt. \\ \hline
legen & liegen & Ich lege mich ins Bett. / Ich liege im Bett. \\ \hline
bringen & tragen & Ich bringe das Buch. / Ich trage das Hemd. \\ \hline
brauchen & müssen & Ich brauche das nicht. / Ich muss das nicht. \\ \hline
werden & sein & Es wird kalt. / Es ist kalt. \\ \hline
\end{tabular}
\end{center}

\chapterend

\part{Anhang}
\pdfbookmark[0]{Teil VII: Anhang}{part7}
\chapter{Anhang: Konjugationstabellen}
\label{cha:Anhang}

Dieser Anhang enthält alle wichtigen Konjugationstabellen als Referenz.

\section{Verbkonjugation: Präsens}
\label{sec:KonjugationPraesens}

\subsection{schwache Verben (regelmäßig)}
\label{sec:SchwacheVerben}

\begin{center}
\begin{tabular}{|l|c|c|c|c|}
\hline
\textbf{} & \textbf{machen} & \textbf{sagen} & \textbf{brauchen} & \textbf{spielen} \\ \hline
ich & mache & sage & brauche & spiele \\ \hline
du & machst & sagst & brauchst & spielst \\ \hline
er/sie/es & macht & sagt & braucht & spielt \\ \hline
wir & machen & sagen & brauchen & spielen \\ \hline
ihr & macht & sagt & braucht & spielt \\ \hline
sie/Sie & machen & sagen & brauchen & spielen \\ \hline
\end{tabular}
\end{center}

\subsection{starke Verben (unregelmäßig)}
\label{sec:StarkeVerben}

\begin{center}
\begin{tabular}{|l|c|c|c|c|}
\hline
\textbf{} & \textbf{lesen} & \textbf{sprechen} & \textbf{geben} & \textbf{nehmen} \\ \hline
ich & lese & spreche & gebe & nehme \\ \hline
du & liest & sprichst & gibst & nimmst \\ \hline
er/sie/es & liest & spricht & gibt & nimmt \\ \hline
wir & lesen & sprechen & geben & nehmen \\ \hline
ihr & lest & sprecht & gebt & nehmt \\ \hline
sie/Sie & lesen & sprechen & geben & nehmen \\ \hline
\end{tabular}
\end{center}

\subsection{haben}
\label{sec:Habentaben}

\begin{center}
\begin{tabular}{|l|c|}
\hline
\textbf{} & \textbf{haben} \\ \hline
ich & habe \\ \hline
du & hast \\ \hline
er/sie/es & hat \\ \hline
wir & haben \\ \hline
ihr & habt \\ \hline
sie/Sie & haben \\ \hline
\end{tabular}
\end{center}

\subsection{sein}
\label{sec:Sein}

\begin{center}
\begin{tabular}{|l|c|}
\hline
\textbf{} & \textbf{sein} \\ \hline
ich & bin \\ \hline
du & bist \\ \hline
er/sie/es & ist \\ \hline
wir & sind \\ \hline
ihr & seid \\ \hline
sie/Sie & sind \\ \hline
\end{tabular}
\end{center}

\subsection{werden}
\label{sec:Werden}

\begin{center}
\begin{tabular}{|l|c|}
\hline
\textbf{} & \textbf{werden} \\ \hline
ich & werde \\ \hline
du & wirst \\ \hline
er/sie/es & wird \\ \hline
wir & werden \\ \hline
ihr & werdet \\ \hline
sie/Sie & werden \\ \hline
\end{tabular}
\end{center}

\section{Verbkonjugation: Präteritum}
\label{sec:KonjugationPraeteritum}

\subsection{schwache Verben}
\label{sec:SchwacheVerbenPraet}

\begin{center}
\begin{tabular}{|l|c|c|c|}
\hline
\textbf{} & \textbf{machen} & \textbf{sagen} & \textbf{brauchen} \\ \hline
ich & machte & sagte & brauchte \\ \hline
du & machtest & sagtest & brauchtest \\ \hline
er/sie/es & machte & sagte & brauchte \\ \hline
wir & machten & sagten & brauchten \\ \hline
ihr & machtet & sagtet & brauchtet \\ \hline
sie/Sie & machten & sagten & brauchten \\ \hline
\end{tabular}
\end{center}

\subsection{starke Verben}
\label{sec:StarkeVerbenPraet}

\begin{center}
\begin{tabular}{|l|c|c|c|c|}
\hline
\textbf{} & \textbf{lesen} & \textbf{sprechen} & \textbf{geben} & \textbf{nehmen} \\ \hline
ich & las & sprach & gab & nahm \\ \hline
du & lasest & sprachst & gabst & nahmst \\ \hline
er/sie/es & las & sprach & gab & nahm \\ \hline
wir & lasen & sprachen & gaben & nahmen \\ \hline
ihr & last & spracht & gabt & nahmt \\ \hline
sie/Sie & lasen & sprachen & gaben & nahmen \\ \hline
\end{tabular}
\end{center}

\section{Modalverben: Präsens}
\label{sec:ModalverbenPraesens}

\begin{center}
\begin{tabular}{|l|c|c|c|c|c|c|}
\hline
\textbf{} & \textbf{dürfen} & \textbf{können} & \textbf{müssen} & \textbf{sollen} & \textbf{wollen} & \textbf{mögen} \\ \hline
ich & darf & kann & muss & soll & will & mag \\ \hline
du & darfst & kannst & musst & sollst & willst & magst \\ \hline
er/sie/es & darf & kann & muss & soll & will & mag \\ \hline
wir & dürfen & können & müssen & sollen & wollen & mögen \\ \hline
ihr & dürft & könnt & müsst & sollt & wollt & mögt \\ \hline
sie/Sie & dürfen & können & müssen & sollen & wollen & mögen \\ \hline
\end{tabular}
\end{center}

\section{Modalverben: Präteritum}
\label{sec:ModalverbenPraeteritum}

\begin{center}
\begin{tabular}{|l|c|c|c|c|c|c|}
\hline
\textbf{} & \textbf{dürfen} & \textbf{können} & \textbf{müssen} & \textbf{sollen} & \textbf{wollen} & \textbf{mögen} \\ \hline
ich & durfte & konnte & musste & sollte & wollte & mochte \\ \hline
du & durftest & konntest & musstest & solltest & wolltest & mochtest \\ \hline
er/sie/es & durfte & konnte & musste & sollte & wollte & mochte \\ \hline
wir & durften & konnten & mussten & sollten & wollten & mochten \\ \hline
ihr & durftet & konntet & musstet & solltet & wolltet & mochtet \\ \hline
sie/Sie & durften & konnten & mussten & sollten & wollten & mochten \\ \hline
\end{tabular}
\end{center}

\section{Modalverben: Konjunktiv II}
\label{sec:ModalverbenKonjunktivII}

\begin{center}
\begin{tabular}{|l|c|c|c|c|c|c|}
\hline
\textbf{} & \textbf{dürfen} & \textbf{können} & \textbf{müssen} & \textbf{sollen} & \textbf{wollen} & \textbf{mögen} \\ \hline
ich & dürfte & könnte & müsste & sollte & wollte & möchte \\ \hline
du & dürftest & könntest & müsstest & solltest & wolltest & möchtest \\ \hline
er/sie/es & dürfte & könnte & müsste & sollte & wollte & möchte \\ \hline
wir & dürften & könnten & müssten & sollten & wollten & möchten \\ \hline
ihr & dürftet & könntet & müsstet & solltet & wolltet & möchtet \\ \hline
sie/Sie & dürften & könnten & müssten & sollten & wollten & möchten \\ \hline
\end{tabular}
\end{center}

\section{Perfekt: Hilfsverben}
\label{sec:PerfektHilfsverben}

\begin{center}
\begin{tabular}{|l|c|c|c|}
\hline
\textbf{} & \textbf{haben} & \textbf{sein} & \textbf{werden} \\ \hline
ich & habe & bin & bin \\ \hline
du & hast & bist & bist \\ \hline
er/sie/es & hat & ist & ist \\ \hline
wir & haben & sind & sind \\ \hline
ihr & habt & seid & seid \\ \hline
sie/Sie & haben & sind & sind \\ \hline
\end{tabular}
\end{center}

\begin{center}
\begin{tabular}{|l|c|c|c|}
\hline
\textbf{Partizip II} & \textbf{gemacht} & \textbf{gegangen} & \textbf{geworden} \\ \hline
\end{tabular}
\end{center}

\section{Konjunktiv I: sein}
\label{sec:KonjunktivIsein}

\begin{center}
\begin{tabular}{|l|c|}
\hline
\textbf{} & \textbf{sein (Konjunktiv I)} \\ \hline
ich & sei \\ \hline
du & seiest \\ \hline
er/sie/es & sei \\ \hline
wir & seien \\ \hline
ihr & seiet \\ \hline
sie/Sie & seien \\ \hline
\end{tabular}
\end{center}

\section{Konjunktiv I: haben}
\label{sec:KonjunktivIhaben}

\begin{center}
\begin{tabular}{|l|c|}
\hline
\textbf{} & \textbf{haben (Konjunktiv I)} \\ \hline
ich & habe \\ \hline
du & habest \\ \hline
er/sie/es & habe \\ \hline
wir & haben \\ \hline
ihr & habet \\ \hline
sie/Sie & haben \\ \hline
\end{tabular}
\end{center}

\section{Konjunktiv I: werden}
\label{sec:KonjunktivIwerden}

\begin{center}
\begin{tabular}{|l|c|}
\hline
\textbf{} & \textbf{werden (Konjunktiv I)} \\ \hline
ich & werde \\ \hline
du & werdest \\ \hline
er/sie/es & werde \\ \hline
wir & werden \\ \hline
ihr & werdet \\ \hline
sie/Sie & werden \\ \hline
\end{tabular}
\end{center}

\section{Konjunktiv II: wäre}
\label{sec:KonjunktivIIwaere}

\begin{center}
\begin{tabular}{|l|c|}
\hline
\textbf{} & \textbf{sein (Konjunktiv II)} \\ \hline
ich & wäre \\ \hline
du & wärest \\ \hline
er/sie/es & wäre \\ \hline
wir & wären \\ \hline
ihr & wärt \\ \hline
sie/Sie & wären \\ \hline
\end{tabular}
\end{center}

\section{Konjunktiv II: hätte}
\label{sec:KonjunktivIIhaette}

\begin{center}
\begin{tabular}{|l|c|}
\hline
\textbf{} & \textbf{haben (Konjunktiv II)} \\ \hline
ich & hätte \\ \hline
du & hättest \\ \hline
er/sie/es & hätte \\ \hline
wir & hätten \\ \hline
ihr & hättet \\ \hline
sie/Sie & hätten \\ \hline
\end{tabular}
\end{center}

\section{ Imperativ}
\label{sec:Imperativ}

\begin{center}
\begin{tabular}{|l|c|c|c|}
\hline
\textbf{} & \textbf{machen} & \textbf{lesen} & \textbf{sein} \\ \hline
du & mach! & lies! & sei! \\ \hline
ihr & macht! & lest! & seid! \\ \hline
Sie & machen Sie! & lesen Sie! & seien Sie! \\ \hline
\end{tabular}
\end{center}

\section{Trennbare Verben: Präsens}
\label{sec:TrennbarPraesens}

\begin{center}
\begin{tabular}{|l|c|c|c|c|}
\hline
\textbf{} & \textbf{abfahren} & \textbf{aufmachen} & \textbf{einladen} & \textbf{vorbereiten} \\ \hline
ich & fahre ab & mache auf & lade ein & bereite vor \\ \hline
du & fährst ab & machst auf & lädst ein & bereitest vor \\ \hline
er/sie/es & fährt ab & macht auf & lädt ein & bereitet vor \\ \hline
wir & fahren ab & machen auf & laden ein & bereiten vor \\ \hline
ihr & fahrt ab & macht auf & ladet ein & bereitet vor \\ \hline
sie/Sie & fahren ab & machen auf & laden ein & bereiten vor \\ \hline
\end{tabular}
\end{center}

\section{Untrennbare Verben: Präsens}
\label{sec:UntrennbarPraesens}

\begin{center}
\begin{tabular}{|l|c|c|c|c|}
\hline
\textbf{} & \textbf{erklären} & \textbf{bezahlen} & \textbf{versuchen} & \textbf{empfangen} \\ \hline
ich & erkläre & bezahle & versuche & empfange \\ \hline
du & erklärst & bezahlst & versuchst & empfängst \\ \hline
er/sie/es & erklärt & bezahlt & versucht & empfängt \\ \hline
wir & erklären & bezahlen & versuchen & empfangen \\ \hline
ihr & erklärt & bezahlt & versucht & empfängt \\ \hline
sie/Sie & erklären & bezahlen & versuchen & empfangen \\ \hline
\end{tabular}
\end{center}

\section{Wichtige unregelmäßige Verben}
\label{sec:WichtigeUnregelmaessige}

\begin{center}
\begin{tabular}{|l|c|c|c|c|}
\hline
\textbf{Infinitiv} & \textbf{Präsens} & \textbf{Präteritum} & \textbf{Perfekt} & \textbf{Bedeutung} \\ \hline
beginnen & beginnt & begann & hat begonnen & beginnen \\ \hline
bleiben & bleibt & blieb & ist geblieben & bleiben \\ \hline
bringen & bringt & brachte & hat gebracht & bringen \\ \hline
denken & denkt & dachte & hat gedacht & denken \\ \hline
fahren & fährt & fuhr & ist gefahren & fahren \\ \hline
finden & findet & fand & hat gefunden & finden \\ \hline
geben & gibt & gab & hat gegeben & geben \\ \hline
gehen & geht & ging & ist gegangen & gehen \\ \hline
gewinnen & gewinnt & gewonnt & hat gewonnen & gewinnen \\ \hline
haben & hat & hatte & hat gehabt & haben \\ \hline
halten & hält & hielt & hat gehalten & halten \\ \hline
hängen & hängt & hing & hat gehangen & hängen \\ \hline
heißen & heißt & hieß & hat geheißen & heißen \\ \hline
kennen & kennt & kannte & hat gekannt & kennen \\ \hline
kommen & kommt & kam & ist gekommen & kommen \\ \hline
können & kann & konnte & hat gekonnt & können \\ \hline
lesen & liest & las & hat gelesen & lesen \\ \hline
liegen & liegt & lag & hat gelegen & liegen \\ \hline
machen & macht & machte & hat gemacht & machen \\ \hline
mögen & mag & mochte & hat gemocht & mögen \\ \hline
müssen & muss & musste & hat gemusst & müssen \\ \hline
nehmen & nimmt & nahm & hat genommen & nehmen \\ \hline
rufen & ruft & rief & hat gerufen & rufen \\ \hline
sagen & sagt & sagte & hat gesagt & sagen \\ \hline
schlafen & schläft & schlief & hat geschlafen & schlafen \\ \hline
schlagen & schlägt & schlug & hat geschlagen & schlagen \\ \hline
schreiben & schreibt & schrieb & hat geschrieben & schreiben \\ \hline
schwimmen & schwimmt & schwamm & ist geschwommen & schwimmen \\ \hline
sehen & sieht & sah & hat gesehen & sehen \\ \hline
sein & ist & war & ist gewesen & sein \\ \hline
senden & sendet & sandte & hat gesandt & senden \\ \hline
singen & singt & sang & hat gesungen & singen \\ \hline
sitzen & sitzt & saß & hat gesessen & sitzen \\ \hline
sollen & soll & sollte & hat gesollt & sollen \\ \hline
sprechen & spricht & sprach & hat gesprochen & sprechen \\ \hline
stehen & steht & stand & hat gestanden & stehen \\ \hline
stimmen & stimmt & stimmte & hat gestimmt & stimmen \\ \hline
tragen & trägt & trug & hat getragen & tragen \\ \hline
treffen & trifft & traf & hat getroffen & treffen \\ \hline
trinken & trinkt & trank & hat getrunken & trinken \\ \hline
tun & tut & tat & hat getan & tun \\ \hline
verstehen & versteht & verstand & hat verstanden & verstehen \\ \hline
werden & wird & wurde & ist geworden & werden \\ \hline
wissen & weiß & wusste & hat gewusst & wissen \\ \hline
wollen & will & wollte & hat gewollt & wollen \\ \hline
ziehen & zieht & zog & hat gezogen & ziehen \\ \hline
\end{tabular}
\end{center}

\section{Artikel-Deklination}
\label{sec:ArtikelDeklination}

\subsection{Bestimmt}
\begin{center}
\begin{tabular}{|l|c|c|c|c|}
\hline
\textbf{Kasus} & \textbf{Maskulin} & \textbf{Feminin} & \textbf{Neutral} & \textbf{Plural} \\ \hline
Nominativ & der & die & das & die \\ \hline
Akkusativ & den & die & das & die \\ \hline
Dativ & dem & der & dem & den \\ \hline
Genitiv & des & der & des & der \\ \hline
\end{tabular}
\end{center}

\subsection{Unbestimmt}
\begin{center}
\begin{tabular}{|l|c|c|c|c|}
\hline
\textbf{Kasus} & \textbf{Maskulin} & \textbf{Feminin} & \textbf{Neutral} & \textbf{Plural} \\ \hline
Nominativ & ein & eine & ein & keine \\ \hline
Akkusativ & einen & eine & ein & keine \\ \hline
Dativ & einem & einer & einem & keinen \\ \hline
Genitiv & eines & einer & eines & keiner \\ \hline
\end{tabular}
\end{center}

\section{Adjektiv-Deklination}
\label{sec:AdjektivDeklination}

\subsection{Mit bestimmtem Artikel}
\begin{center}
\begin{tabular}{|l|c|c|c|c|}
\hline
\textbf{Kasus} & \textbf{Maskulin} & \textbf{Feminin} & \textbf{Neutral} & \textbf{Plural} \\ \hline
Nominativ & der gute & die gute & das gute & die guten \\ \hline
Akkusativ & den guten & die gute & das gute & die guten \\ \hline
Dativ & dem guten & der guten & dem guten & den guten \\ \hline
Genitiv & des guten & der guten & des guten & der guten \\ \hline
\end{tabular}
\end{center}

\subsection{Mit unbestimmtem Artikel}
\begin{center}
\begin{tabular}{|l|c|c|c|c|}
\hline
\textbf{Kasus} & \textbf{Maskulin} & \textbf{Feminin} & \textbf{Neutral} & \textbf{Plural} \\ \hline
Nominativ & ein guter & eine gute & ein gutes & keine guten \\ \hline
Akkusativ & einen guten & eine gute & ein gutes & keine guten \\ \hline
Dativ & einem guten & einer guten & einem guten & keinen guten \\ \hline
Genitiv & eines guten & einer guten & eines guten & keiner guten \\ \hline
\end{tabular}
\end{center}

\section{Relativpronomen (Anhang)}
\label{sec:RelativpronomenAnhang}

\begin{center}
\begin{tabular}{|l|c|c|c|c|}
\hline
\textbf{Kasus} & \textbf{Maskulin} & \textbf{Feminin} & \textbf{Neutral} & \textbf{Plural} \\ \hline
Nominativ & der & die & das & die \\ \hline
Akkusativ & den & die & das & die \\ \hline
Dativ & dem & der & dem & denen \\ \hline
Genitiv & dessen & deren & dessen & deren \\ \hline
\end{tabular}
\end{center}

\section{Fragepronomen}
\label{sec:Fragepronomen}

\begin{center}
\begin{tabular}{|l|c|c|c|c|}
\hline
\textbf{Kasus} & \textbf{Maskulin} & \textbf{Feminin} & \textbf{Neutral} & \textbf{Plural} \\ \hline
Nominativ & wer? & was? \\ \hline
Akkusativ & wen? & was? \\ \hline
Dativ & wem? \\ \hline
Genitiv & wessen? \\ \hline
\end{tabular}
\end{center}

\section{Wechselpräpositionen}
\label{sec:Wechselpraepositionen}

\begin{center}
\begin{tabular}{|l|c|c|}
\hline
\textbf{Präposition} & \textbf{Wo? (Dativ)} & \textbf{Wohin? (Akkusativ)} \\ \hline
in & in der Schule & in die Schule \\ \hline
an & an der Wand & an die Wand \\ \hline
auf & auf dem Tisch & auf den Tisch \\ \hline
unter & unter dem Tisch & unter den Tisch \\ \hline
vor & vor dem Haus & vor das Haus \\ \hline
hinter & hinter dem Haus & hinter das Haus \\ \hline
neben & neben dem Haus & neben das Haus \\ \hline
zwischen & zwischen den Bäumen & zwischen die Bäume \\ \hline
über & über dem Tisch & über den Tisch \\ \hline
\end{tabular}
\end{center}

\section{Übersicht: Präpositionen nach Kasus}
\label{sec:PraepositionenNachKasus}

\subsection{Dativ-Präpositionen}
\begin{center}
\begin{tabular}{|l|l|}
\hline
\textbf{Präposition} & \textbf{Bedeutung/Beispiel} \\ \hline
aus & Herkunft: Ich komme aus Berlin. \\ \hline
bei & Nähe: Ich bin bei Maria. \\ \hline
mit & Begleitung: Ich fahre mit dem Bus. \\ \hline
nach & Richtung: Nach dem Essen... \\ \hline
seit & Zeit: Seit zwei Jahren... \\ \hline
von & Herkunft: Das Buch ist von ihm. \\ \hline
zu & Richtung: Ich gehe zur Schule. \\ \hline
ab & Zeitpunkt: Ab morgen... \\ \hline
außer & Ausnahme: Außer mir war niemand da. \\ \hline
\end{tabular}
\end{center}

\subsection{Akkusativ-Präpositionen}
\begin{center}
\begin{tabular}{|l|l|}
\hline
\textbf{Präposition} & \textbf{Bedeutung/Beispiel} \\ \hline
bis & Richtung/Zeit: Bis morgen! \\ \hline
durch & Richtung: Durch den Wald \\ \hline
gegen & Richtung: Gegen die Wand \\ \hline
ohne & Ausschluss: Ohne dich \\ \hline
um & Richtung: Um die Ecke \\ \hline
\end{tabular}
\end{center}

\subsection{Genitiv-Präpositionen}
\begin{center}
\begin{tabular}{|l|l|}
\hline
\textbf{Präposition} & \textbf{Bedeutung/Beispiel} \\ \hline
trotz & trotz des Regens... \\ \hline
wegen & Wegen des Verkehrs... \\ \hline
während & Während des Films... \\ \hline
statt & Statt des Bruders... \\ \hline
außerhalb & Außerhalb der Stadt \\ \hline
innerhalb & Innerhalb der Woche \\ \hline
aufgrund & Aufgrund der Situation... \\ \hline
oberhalb & Oberhalb des Berges \\ \hline
unterhalb & Unterhalb des Tales \\ \hline
diesseits & Diesseits des Flusses \\ \hline
jenseits & Jenseits des Meeres \\ \hline
\end{tabular}
\end{center}

\section{Reflexive Verben mit Präpositionen (sich + Präposition)}
\label{sec:ReflexivPraep}

Viele reflexive Verben werden mit festen Präpositionen verwendet:

\subsection{sich an... gewöhnen}
\begin{center}
\begin{tabular}{|l|c|c|}
\hline
\textbf{Pronomen} & \textbf{Beispiel} & \textbf{Bedeutung} \\ \hline
ich & Ich gewöhne mich an das Wetter. & adapt to \\ \hline
du & Du gewöhnst dich an die neue Schule. & adapt to \\ \hline
er/sie/es & Er gewöhnt sich an den Job. & adapt to \\ \hline
wir & Wir gewöhnen uns an die Stadt. & adapt to \\ \hline
\end{tabular}
\end{center}

\subsection{sich auf... freuen}
\begin{center}
\begin{tabular}{|l|c|c|}
\hline
\textbf{Pronomen} & \textbf{Beispiel} & \textbf{Bedeutung} \\ \hline
ich & Ich freue mich auf den Urlaub. & look forward to \\ \hline
du & Du freust dich auf das Fest. & look forward to \\ \hline
er/sie/es & Sie freut sich auf dich. & look forward to \\ \hline
wir & Wir freuen uns auf Weihnachten. & look forward to \\ \hline
\end{tabular}
\end{center}

\subsection{sich auf... vorbereiten}
\begin{center}
\begin{tabular}{|l|c|c|}
\hline
\textbf{Pronomen} & \textbf{Beispiel} & \textbf{Bedeutung} \\ \hline
ich & Ich bereite mich auf die Prüfung vor. & prepare for \\ \hline
du & Du bereitest dich auf das Interview vor. & prepare for \\ \hline
er/sie/es & Er bereitet sich auf die Reise vor. & prepare for \\ \hline
wir & Wir bereiten uns auf die Präsentation vor. & prepare for \\ \hline
\end{tabular}
\end{center}

\subsection{sich über... ärgern}
\begin{center}
\begin{tabular}{|l|c|c|}
\hline
\textbf{Pronomen} & \textbf{Beispiel} & \textbf{Bedeutung} \\ \hline
ich & Ich ärgere mich über den Lärm. & be annoyed at \\ \hline
du & Du ärgerst dich über die Note. & be annoyed at \\ \hline
er/sie/es & Er ärgert sich über den Chef. & be annoyed at \\ \hline
wir & Wir ärgern uns über das Wetter. & be annoyed at \\ \hline
\end{tabular}
\end{center}

\subsection{sich über... informieren}
\begin{center}
\begin{tabular}{|l|c|c|}
\hline
\textbf{Pronomen} & \textbf{Beispiel} & \textbf{Bedeutung} \\ \hline
ich & Ich informiere mich über die Angebote. & inform about \\ \hline
du & Du informierst dich über das Programm. & inform about \\ \hline
er/sie/es & Er informiert sich über die Preise. & inform about \\ \hline
wir & Wir informieren uns über die Bedingungen. & inform about \\ \hline
\end{tabular}
\end{center}

\section{Nomen: Artikel und Plural}
\label{sec:NomenArtikel}

\subsection{Maskulin (der)}
\begin{center}
\begin{tabular}{|l|c|c|l|}
\hline
\textbf{Nomen} & \textbf{Artikel} & \textbf{Plural} & \textbf{Bedeutung} \\ \hline
der Mann & der & die Männer & man \\ \hline
der Junge & der & die Jungen & boy \\ \hline
der Vater & der & die Väter & father \\ \hline
der Sohn & der & die Söhne & son \\ \hline
der Bruder & der & die Brüder & brother \\ \hline
der Freund & der & die Freunde & friend \\ \hline
der Lehrer & der & die Lehrer & teacher \\ \hline
der Arzt & der & die Ärzte & doctor \\ \hline
der Student & der & die Studenten & student \\ \hline
der Mensch & der & die Menschen & human \\ \hline
der Herr & der & die Herren & gentleman \\ \hline
der Kunde & der & die Kunden & customer \\ \hline
der Kollege & der & die Kollegen & colleague \\ \hline
der Chef & der & die Chefs & boss \\ \hline
der Hund & der & die Hunde & dog \\ \hline
der Vogel & der & die Vögel & bird \\ \hline
der Tisch & der & die Tische & table \\ \hline
der Stuhl & der & die Stühle & chair \\ \hline
der Computer & der & die Computer & computer \\ \hline
der Tag & der & die Tage & day \\ \hline
\end{tabular}
\end{center}

\subsection{Feminin (die)}
\begin{center}
\begin{tabular}{|l|c|c|l|}
\hline
\textbf{Nomen} & \textbf{Artikel} & \textbf{Plural} & \textbf{Bedeutung} \\ \hline
die Frau & die & die Frauen & woman \\ \hline
die Mutter & die & die Mütter & mother \\ \hline
die Schwester & die & die Schwestern & sister \\ \hline
die Tochter & die & die Töchter & daughter \\ \hline
die Freundin & die & die Freundinnen & (female) friend \\ \hline
die Lehrerin & die & die Lehrerinnen & (female) teacher \\ \hline
die Ärztin & die & die Ärztinnen & (female) doctor \\ \hline
die Studentin & die & die Studentinnen & (female) student \\ \hline
die Katze & die & die Katzen & cat \\ \hline
die Blume & die & die Blumen & flower \\ \hline
die Schule & die & die Schulen & school \\ \hline
die Stadt & die & die Städte & city \\ \hline
die Uhr & die & die Uhren & clock/watch \\ \hline
die Lampe & die & die Lampen & lamp \\ \hline
die Tür & die & die Türen & door \\ \hline
die Hand & die & die Hände & hand \\ \hline
die Zeit & die & die Zeiten & time \\ \hline
die Arbeit & die & die Arbeiten & work \\ \hline
die Frage & die & die Fragen & question \\ \hline
die Sprache & die & die Sprachen & language \\ \hline
\end{tabular}
\end{center}

\subsection{Neutral (das)}
\begin{center}
\begin{tabular}{|l|c|c|l|}
\hline
\textbf{Nomen} & \textbf{Artikel} & \textbf{Plural} & \textbf{Bedeutung} \\ \hline
das Kind & das & die Kinder & child \\ \hline
das Buch & das & die Bücher & book \\ \hline
das Haus & das & die Häuser & house \\ \hline
das Auto & das & die Autos & car \\ \hline
das Bild & das & die Bilder & picture \\ \hline
das Wort & das & die Wörter & word \\ \hline
das Jahr & das & die Jahre & year \\ \hline
das Zimmer & das & die Zimmer & room \\ \hline
das Fenster & das & die Fenster & window \\ \hline
das Telefon & das & die Telefone & telephone \\ \hline
das Radio & das & die Radios & radio \\ \hline
das Foto & das & die Fotos & photo \\ \hline
das Baby & das & die Babys & baby \\ \hline
das Problem & das & die Probleme & problem \\ \hline
das Land & das & die Länder & country/land \\ \hline
das Meer & das & die Meere & sea \\ \hline
das Tier & das & die Tiere & animal \\ \hline
das Wasser & das & - & water \\ \hline
das Brot & das & - & bread \\ \hline
das Geld & das & - & money \\ \hline
\end{tabular}
\end{center}

\section{Nützliche Sätze für Meinungen und Optionen}
\label{sec:Meinungen}

\subsection{Meinungen ausdrücken}
\begin{center}
\begin{tabular}{|l|p{8cm}|}
\hline
\textbf{Satz} & \textbf{Bedeutung} \\ \hline
Ich finde, dass... & I think that... \\ \hline
Ich meine, dass... & I believe that... \\ \hline
Ich bin der Meinung, dass... & I am of the opinion that... \\ \hline
Meiner Meinung nach... & In my opinion... \\ \hline
Ich denke, dass... & I think that... \\ \hline
Ich halte es für... & I consider it to be... \\ \hline
Es scheint mir, dass... & It seems to me that... \\ \hline
Ich bin überzeugt, dass... & I am convinced that... \\ \hline
Es ist klar, dass... & It is clear that... \\ \hline
Natürlich... & Of course... \\ \hline
Klar... & Obviously... \\ \hline
Leider... & Unfortunately... \\ \hline
Ich stimme zu. & I agree. \\ \hline
Ich stimme nicht zu. & I disagree. \\ \hline
Das sehe ich anders. & I see it differently. \\ \hline
Das ist richtig. & That is correct. \\ \hline
Das ist falsch. & That is wrong. \\ \hline
\end{tabular}
\end{center}

\subsection{Optionen und Alternativen vorschlagen}
\begin{center}
\begin{tabular}{|l|p{8cm}|}
\hline
\textbf{Satz} & \textbf{Bedeutung} \\ \hline
Einerseits... andererseits... & On one hand... on the other hand... \\ \hline
Es gibt mehrere Möglichkeiten. & There are several possibilities. \\ \hline
Man könnte... & One could... \\ \hline
Man könnte auch... & One could also... \\ \hline
Wie wäre es mit...? & How about...? \\ \hline
Was hältst du von...? & What do you think of...? \\ \hline
Ich schlage vor, dass... & I suggest that... \\ \hline
Ich empfehle... & I recommend... \\ \hline
Es wäre besser, wenn... & It would be better if... \\ \hline
Anstelle von... könnte man... & Instead of... one could... \\ \hline
Alternativ könnte... & Alternatively... could... \\ \hline
Eine andere Möglichkeit wäre... & Another possibility would be... \\ \hline
Das beste wäre... & The best would be... \\ \hline
Im Gegensatz dazu... & In contrast to that... \\ \hline
\end{tabular}
\end{center}

\subsection{Zustimmen und Widersprechen}
\begin{center}
\begin{tabular}{|l|p{8cm}|}
\hline
\textbf{Satz} & \textbf{Bedeutung} \\ \hline
Ja, das stimmt. & Yes, that's right. \\ \hline
Genau so ist es. & Exactly. \\ \hline
Da hast du recht. & You're right there. \\ \hline
Ich bin ganz deiner Meinung. & I am completely of your opinion. \\ \hline
Das sehe ich genauso. & I see it the same way. \\ \hline
Nein, das sehe ich anders. & No, I see it differently. \\ \hline
Da muss ich widersprechen. & I have to disagree there. \\ \hline
Das ist nicht richtig. & That's not correct. \\ \hline
Ich sehe das nicht so. & I don't see it that way. \\ \hline
Das kann man nicht generalisieren. & You can't generalize that. \\ \hline
\end{tabular}
\end{center}

\section{Wichtige Wortverbindungen}
\label{sec:Wortverbindungen}

\subsection{Nomen + Präposition}
\begin{center}
\begin{tabular}{|l|c|l|}
\hline
\textbf{Nomen} & \textbf{+ Präposition} & \textbf{Beispiel} \\ \hline
die Angst & vor & Angst vor dem Hund \\ \hline
der Spaß & an & Spaß am Sport \\ \hline
das Interesse & an & Interesse an Musik \\ \hline
die Hoffnung & auf & Hoffnung auf Erfolg \\ \hline
die Freude & an & Freude am Lernen \\ \hline
der Ärger & über & Ärger über die Entscheidung \\ \hline
die Frage & nach & die Frage nach dem Sinn \\ \hline
der Grund & für & der Grund für die Verspätung \\ \hline
das Problem & mit & das Problem mit dem Computer \\ \hline
die Lösung & für & die Lösung für das Problem \\ \hline
\end{tabular}
\end{center}

\subsection{Adjektiv + Präposition}
\begin{center}
\begin{tabular}{|l|c|l|}
\hline
\textbf{Adjektiv} & \textbf{+ Präposition} & \textbf{Beispiel} \\ \hline
stolz & auf & stolz auf die Leistung \\ \hline
zufrieden & mit & zufrieden mit dem Ergebnis \\ \hline
glücklich & über & glücklich über die Nachricht \\ \hline
traurig & über & traurig über das Ende \\ \hline
wütend & auf & wütend auf den Chef \\ \hline
nervös & wegen & nervös wegen der Prüfung \\ \hline
abhängig & von & abhängig von den Eltern \\ \hline
interessiert & an & interessiert an Kunst \\ \hline
reich & an & reich an Erfahrungen \\ \hline
arm & an & arm an Zeit \\ \hline
\end{tabular}
\end{center}

\chapterend

\end{document}